\documentclass[11pt]{article}
\usepackage[a4paper,margin=27mm]{geometry}
\usepackage{amsmath,amssymb,amsthm,mathtools}
\usepackage[T1]{fontenc}
\usepackage{lmodern}
\usepackage{booktabs,longtable,array,microtype}
\usepackage{xurl}
\usepackage[hidelinks]{hyperref}
\hypersetup{pdftitle={Explicit Stability of the Weak Simplex Bound},
  pdfauthor={Research draft}}
\Urlmuskip=0mu plus 1mu
\newtheorem{theorem}{Theorem}
\newtheorem{lemma}[theorem]{Lemma}
\newtheorem{proposition}[theorem]{Proposition}
\newtheorem{corollary}[theorem]{Corollary}
\theoremstyle{remark}\newtheorem{remark}[theorem]{Remark}
\newcommand{\R}{\mathbb R}
\newcommand{\E}{\mathbb E}
\newcommand{\Pp}{\mathbb P}
\newcommand{\one}{\mathbf 1}
\newcommand{\tr}{\operatorname{tr}}
\newcommand{\diag}{\operatorname{diag}}
\newcommand{\rank}{\operatorname{rank}}
\newcommand{\Id}{I_m}
\newcommand{\norm}[1]{\left\lVert#1\right\rVert}
\newcommand{\Pc}{P_{\mathrm c}}
\title{Explicit Stability of the Weak Simplex Bound}
\author{Research draft}
\date{6 November 2026}
\begin{document}
\maketitle
\begin{abstract}
Equal-energy Gaussian signaling uses the regular simplex as an optimal reference
when the signal dimension permits it. Quantifying geometric recovery from a
small decoding deficit requires a bound that remains valid near singular
Gram matrices. We give explicit global quadratic bounds for Gram and vertex
distance at every fixed positive channel parameter, together with a separate
centroid penalty. A configuration-metric expansion separates two in-plane
sectors from a transverse centroid sector and identifies the local least-loss
directions. We prove unique ordered channel-parameter crossings and identify
the exact minimizers on all sufficiently small Procrustes shells in the
strict-normal regime, uniformly in the ambient dimension. Their minimal
deficit has an explicit quartic expansion with a sixth-order remainder.
The global bounds, mean-width stability, transition and small-shell
results, including the remainder, are formalized in Lean. For four- and 16-signal constellations, directional coefficients
predict central perturbation losses more accurately in aggregate than two
scalar controls across 108 tested cells. A separate 20-run learning comparison
attributes improved matched-decoder performance to ordinary simplex Gram
regularization, rather than to a new optimization method. The evaluated global
recovery constants range from approximately $10^{-39}$ to $10^{-234}$ and
give no practical finite-precision geometry certificate. The results provide
an explicit mathematical rate and a directional local diagnostic, while
practical certification requires sharper constants and effective numerical
bounds for the local radius and remainder constant.
\end{abstract}

\section{Setting and the quantitative question}
Fix $m\ge2$, put $n=m-1$, $\alpha=n/m$, $J=\one\one^{\mathsf T}$ and
$P=I_m-J/m$. The regular unit simplex has Gram matrix $G_\Delta=P/\alpha$.
For a correlation matrix $G$ (positive semidefinite with unit diagonal), define
\begin{equation}\label{eq:normalization}
 R=\alpha G+J/m,\qquad
 \mathcal A_m=\{R:\diag R=\one,\ R-J/m\succeq0\}.
\end{equation}
Thus $G=G_\Delta$ corresponds to $R=I_m$. Write
\[
 F_R(c)=\Pp\{X_i\le c\text{ for all }i\},\qquad X\sim N(0,R),
 \quad D_R(c)=F_R(c)-\Phi(c)^m,
\]
where $\Phi$ and $\phi$ are the standard normal distribution and density.
We use the comparison $D_R(c)\ge0$ proved in \cite{Mulgund}, as well as its
centered product inequality and variational construction. Pointwise strictness was established independently by Su et al.\ \cite{Su}
and by Mulgund, whose v4 explicitly records this parallel development
\cite[Remark 2.3 and Appendix B]{MulgundLatest}. Su et al.\ also explicitly
suggest a quantitative deficit estimate. The present argument replaces a
strict positive gap by a computable lower bound. We cite v2 for the particular
product and variational statements used in the proof, not as a description
of the current state of Mulgund's work.

For unit signals $x_1,\ldots,x_m$ in $\R^d$, let $G_{ij}=\langle x_i,x_j\rangle$
and consider $Y=\lambda x_I+Z$, with $I$ uniform, $Z\sim N(0,I_d)$, and $\lambda>0$.
Maximum-likelihood decoding gives the identity
\begin{equation}\label{eq:pc}
 \Pc(G;\lambda)=\frac{e^{-\lambda^2/2}}m
       \E e^{\lambda\max_i\xi_i},\qquad \xi\sim N(0,G).
\end{equation}
All deficits below compare to the unit centered regular simplex at the same
operating parameter. The integer $m$ is the number of signals, and $d$ is the
ambient dimension of the signal space. Decoding constants and unnormalized
expected-maximum constants depend on $m$ and, for decoding, on $\lambda$, but
not on $d$. Only those constants are called uniform in the signal-space
dimension. The normalized mean-width lower bound \eqref{eq:width-stability}
retains the factor $2/\E\norm g$ for $g\sim N(0,I_d)$, and that prefactor
depends on $d$. No positive lower constant uniform in $m$, and no positive
lower constant uniform over all positive channel parameters, is asserted.

Local optimality of simplex codes is classical.
Litvak attributes fixed-parameter local maximality to Balakrishnan
\cite[Theorem 2.5]{Litvak}, whose original article is \cite{Balakrishnan}.
Weber already expresses the general second variation through repeated,
adjacent and disjoint edge pairs \cite[Lemma 5.8, equation (5.102)]{Weber}.
His three quadratic invariants are the squared edge norm, squared total row
sum and sum of squared row sums. On his feasible zero-total-sum face,
the row sums vanish and (5.105) gives the restricted form used in Theorem 5.5.
Equations (5.107)--(5.108) integrate the distribution-function variation into
the decoding functional. Thus edge-type decomposition and local simplex
optimality are established antecedents of the calculation below.
We claim no priority for these local ideas or for the existence of a Hessian.
Our global quantitative contribution is the specified rate and its
conversion into configuration distance. The local formulas supply explicit
normalization and feasible-sequence criteria for using that rate.
Section~\ref{sec:new-results} additionally proves unique channel transitions
and an exact, rather than only second-order, optimizer on small shells in
the strict-normal regime. Its complete analytic and geometric proofs are
in Appendices~\ref{ap:transitions} and~\ref{ap:normal}; the precise scope of
the Lean verification is given in Appendix~\ref{ap:lean}.

Mean-width stability under John/L\"owner ellipsoid constraints \cite{BFH}
concerns a different admissible class. The unsmoothed comparison in
\cite[Theorem 1.1]{MulgundFull} dominates Gaussian maxima at every real
threshold and characterizes equality at positive thresholds.
Its equality statement does not extend to every nonpositive threshold.
We use the cited comparison and centered product inequality as inputs.
Quantitative Gaussian correlation inequalities also predate this work
\cite{DNS}. Our rectangle estimate is a specialized input.

Friedland, Gu\'edon, and Souli \cite{FGS} prove an explicit stability theorem
for a different functional. For $\beta\in\R$ they set
$Z_\beta(X)=\sum_{i=1}^N e^{\beta X_i}$, where $X$ is centered Gaussian with
unit variances and correlation matrix $C\in E_N$. The set $E_N$ consists of
all $N\times N$ correlation matrices, and $\Delta_N$ is the regular-simplex
correlation. That matrix equals $G_\Delta$. The symbol $\beta$ in this
discussion is theirs, not the scalar in \eqref{eq:bounds}. Their coefficients
$\mu_{N,q,\beta}$ and $\eta_{N,q,\beta}$ are not the constants $\mu$ and
$\eta$ in \eqref{eq:codeconstants} and \eqref{eq:eta}.

Their Theorem~1.1 is a phase transition of generalized $L_q$-means, for every
real $q$, every real $\beta$, and every $C\in E_N$. The simplex correlation
maximizes $\|Z_\beta(X)\|_q$ when $q<1$. The mean equals $N e^{\beta^2/2}$
for every such $C$ when $q=1$. The simplex correlation minimizes the mean
when $q>1$. The quantitative statement is narrower. Theorem~1.3 requires
$q>1$ and $\beta\neq 0$, and concludes
\[
 \E[Z_\beta(X)^q]-\E[Z_\beta(S)^q]
 \ge \eta_{N,q,\beta}\,\one^{\mathsf T}C\one
  +\frac{\mu_{N,q,\beta}}{2}\norm{C-\Delta_N}_F^2,
\]
where $S\sim N(0,\Delta_N)$ and both coefficients are positive. The summand
$\eta_{N,q,\beta}\,\one^{\mathsf T}C\one$ measures centroid defect, and the
Frobenius summand measures shape defect. The
same modulus $\mu_{N,q,\beta}$, their (15), is the strong-convexity modulus
of $C\mapsto\E[Z_\beta(X)^q]$ on all of $E_N$ in their Theorem~1.2. Their
(15) and (16) vanish at $\beta=0$, so both the modulus and the two stability
terms are zero in that case. Theorem~1.1 remains valid at $\beta=0$, where
$Z_0$ is the constant $N$.

The estimates below are not those statements. Since
$\alpha G_\Delta=P=I_m-J/m$, the normalization \eqref{eq:normalization} gives
$R-I_m=\alpha(G-G_\Delta)$ and identifies $\mathcal A_m$ with the image of
$E_m$. Since $\Delta_m=G_\Delta$, it follows that
$\norm{R-I_m}_F=\alpha\norm{G-\Delta_m}_F$. The constant
$\kappa_{m,[u,v]}$ depends on a compact threshold interval and is not
$\mu_{N,q,\beta}$. Theorem~1.3 controls moments of order $q>1$ of $Z_\beta$
on $E_N$. The deficit $D_R(c)$ is an orthant probability on $\mathcal A_m$,
uniform for $c$ in a compact interval, and the bound includes singular $R$.
The inequality $D_R(c)\ge 0$ has the direction of their Corollary~1.5 applied
to the indicator of $(-\infty,c]$. That corollary assumes a correlation
matrix $R$ with $R-J/N\succeq 0$, which is the class $\mathcal A_N$ with
their $N$ equal to the number of signals $m$. Remark~1.6 records the
resulting non-strict bound as generalizing an orthant inequality of Mulgund.
The corollary does not give a quadratic rate. For a Gram matrix of unit
vectors, their centroid summand and Lemma~\ref{lem:centroid} both control
$\norm{\sum_i x_i}^2=\one^{\mathsf T}G\one$, through different functionals
and different coefficients. Decoding and mean-width stability below concern
$\E e^{\lambda\max_i\xi_i}$ and $\E\max_i\xi_i$. They record that
$\beta^{-1}\E\log Z_\beta$ tends to $\E\max_i X_i$ as $\beta\to+\infty$. At
$q=0$, Theorem~1.1 compares $\E\log Z_\beta$ at fixed $\beta$. Theorem~1.3
covers neither statement, because it assumes $q>1$. The full-class local
liminf, including the split between $m=3$ and $m\ge 4$, is not a conclusion
of Theorem~1.1 or Theorem~1.3.

\begin{proposition}[What compactness supplies]\label{prop:compactness}
Let $\mathcal K$ be a compact metric space, $Z\subset\mathcal K$ a nonempty
closed set, and $f:\mathcal K\to[0,\infty)$ continuous with zero set exactly
$Z$. Then $f(x_j)\to0$ implies $d(x_j,Z)\to0$. If, in addition,
$f(x)\ge a\,d(x,Z)^2$ whenever $d(x,Z)<\rho$, with $a,\rho>0$,
then there exists $b>0$ such that $f(x)\ge b\,d(x,Z)^2$ on $\mathcal K$.
\end{proposition}
\begin{proof}
For each $\epsilon>0$, compactness gives a positive minimum of $f$ on
$\{d(x,Z)\ge\epsilon\}$ whenever that set is nonempty. This proves the
first assertion. For the second, let $M=\max_{\mathcal K}d(x,Z)$.
If $M=0$ or the set $\{d(x,Z)\ge\rho\}$ is empty, the assertion is
immediate. Otherwise its minimum $\eta$ is positive, and
$b=\min\{a,\eta/M^2\}$ works.
\end{proof}
For fixed $m$, the correlation-matrix domain is compact. The analogous
vertex formulation uses a fixed ambient space and the closed simplex orbit.
Thus a nonexplicit global quadratic bound is a standard consequence of
uniqueness and a local quadratic bound in the \emph{same} metric.
Uniqueness and compactness alone give no power law. Nor are the Gram and
vertex metrics locally equivalent at the rank-deficient simplex, as
Section~\ref{sec:spectrum} shows through an explicit transverse example.
These distinctions delimit the role of our global proof and constants.

\section{Explicit constants and the global comparison}
Define the scalar functions
\begin{equation}\label{eq:scalar}
 r(s)=\frac{\phi(s)}{\Phi(s)},\qquad H(s)=s+r(s),\qquad
 \Psi(H(s))=\log\Phi(s)+\frac12r(s)^2.
\end{equation}
The map $H:\R\to(0,\infty)$ is an increasing bijection.
Moreover $\Psi:(0,\infty)\to(-\infty,0)$ is an increasing, strictly concave
bijection, with
\begin{equation}\label{eq:psi-derivative}
 \Psi'(H(s))=r(s),\quad
 \Psi''(H(s))=-\frac{r(s)H(s)}{1-r(s)H(s)}<0.
\end{equation}
The denominator is the variance of a standard normal conditioned on $Z\le s$.
These scalar identities also fix the meaning of the one-dimensional inverses
and maximizations used in our constants.

Fix a finite interval $[u,v]$, allowing $u=v$, and set
\begin{align}
 L_0&=m\log\Phi(u),&
 \beta&=\Psi^{-1}(L_0),& A_0&=\Psi'(\beta),\notag\\
 B_0&=v+\sqrt{-2L_0/m}+\sqrt{-2\alpha L_0},&
 \ell&=\beta/\sqrt2,& T&=\sqrt2 B_0+\ell.\label{eq:bounds}
\end{align}
All these quantities are finite, and $\beta,\ell,B_0,T,A_0$ are positive.
For $t>0$ put $b(t)=2\Phi(t)-1$. Define
\begin{equation}\label{eq:K}
 K_{m,[u,v]}=
 \frac{e^{-mA_0\beta}\ell^2e^{-4T^2/3}b(\ell)^{2m-2}}
      {4\pi m(m-1)},\qquad
 \kappa_{m,[u,v]}=
 \frac{\Phi(u)^m K_{m,[u,v]}}
      {1+\sqrt{1+K_{m,[u,v]}m(m-1)}}.
\end{equation}

\begin{theorem}[Global quantitative comparison]\label{thm:global}
For every $R\in\mathcal A_m$ and $c\in[u,v]$,
\begin{equation}\label{eq:global}
 D_R(c)\ge\Phi(u)^m\left(\sqrt{1+K_{m,[u,v]}\norm{R-I_m}_F^2}-1\right)
 \ge\kappa_{m,[u,v]}\norm{R-I_m}_F^2.
\end{equation}
The conclusion includes singular $R$.
\end{theorem}

We first quantify the two-coordinate input. Its conversion to a full rectangle
uses the Gaussian correlation inequality \cite{Royen}.
\begin{lemma}[Quantitative rectangle gap]\label{lem:rectangle}
Let $V\sim N(0,S)$, where $S$ is any $m\times m$ correlation matrix, possibly
singular. If $t_i\in[\ell,T]$, with $0<\ell\le T<\infty$, then
\begin{equation}\label{eq:rectangle}
 \Pp\{|V_i|\le t_i\ \forall i\}-\prod_i b(t_i)
 \ge \frac{\ell^2e^{-4T^2/3}b(\ell)^{m-2}}{4\pi m(m-1)}
       \norm{S-I_m}_F^2.
\end{equation}
\end{lemma}
\begin{proof}
For a bivariate standard normal with correlation $\rho$, let
$B_\rho(a,b)=\Pp\{|V_1|\le a,|V_2|\le b\}$.
This is even in $\rho$. Differentiating the Gaussian integral (equivalently,
integrating $\partial_\rho f_\rho=\partial_x\partial_y f_\rho$) gives, for $0<\rho<1$,
\[
 \partial_\rho B_\rho(a,b)=
 \frac{2}{\pi\sqrt{1-\rho^2}}
 e^{-(a^2+b^2)/(2(1-\rho^2))}
 \sinh\!\left(\frac{\rho ab}{1-\rho^2}\right).
\]
For $a,b\in[\ell,T]$ and $0\le\rho\le1/2$, this derivative is at least
$(2\ell^2/\pi)e^{-4T^2/3}\rho$. It is nonnegative on $[0,1)$, so integration gives
\[
 B_\rho(a,b)-b(a)b(b)
 \ge \frac{\ell^2}{\pi}e^{-4T^2/3}\min\{|\rho|,1/2\}^2
 \ge \frac{\ell^2}{4\pi}e^{-4T^2/3}\rho^2.
\]
Continuity extends this non-strict quantitative inequality to $|\rho|=1$;
the rectangle boundary has zero probability since the marginals have no atoms.
Choose $i<j$ with maximal $S_{ij}^2$. Group this two-coordinate rectangle
with the remaining symmetric slabs using the Gaussian correlation inequality.
Its deficit is multiplied by at least $b(\ell)^{m-2}$.
Finally $\max_{i<j}S_{ij}^2\ge\norm{S-I_m}_F^2/[m(m-1)]$.
For singular $S$, apply the correlation inequality to the linear preimages of
these sets under a representation $V=TW$ with nondegenerate $W$.
\end{proof}

\begin{proof}[Proof of Theorem~\ref{thm:global}]
Represent $G=(R-J/m)/\alpha$ by unit vectors $x_i$. We recall precisely the
two inputs from \cite{Mulgund}. The centered product inequality says
\begin{equation}\label{eq:product}
 \E\prod_i f_i(Y_i)\ge\prod_i\int f_i\,d\gamma_1
\end{equation}
for any Gaussian correlation matrix of $Y$, and bounded nonnegative
log-concave factors with positive Gaussian integrals and zero Gaussian first
moments. The adaptive tilts arise from maximizing
\[
 \mathcal J_c(q,w)=\sum_i\Psi(q+\alpha\langle x_i,w\rangle)
             -\frac\alpha2\norm w^2-\frac m2(q-c)^2
\]
on the domain where all scalar arguments are positive. Its maximizer has
\begin{equation}\label{eq:tilt}
b_i=H(s_i)=c+(Ra)_i,\quad a_i=r(s_i),\qquad
 \mathcal J_*\ge m\log\Phi(c)\ge L_0.
\end{equation}
For clarity, stationarity of the strictly concave variational functional gives
$w_*=\sum_i a_i x_i$ and $q_*-c=m^{-1}\sum_i a_i$, and hence
\[
 \mathcal J_*=\sum_i\Psi(b_i)-\tfrac12 a^{\mathsf T}Ra.
\]
The maximizer exists in the positive-argument domain: the quadratic terms
control infinity, whereas $\Psi(y)\to-\infty$ as $y\downarrow0$.
We now extract bounds not requiring any inverse of $R$.
Since every summand $\Psi(b_i)$ and both negative quadratic terms are
nonpositive, $\Psi(b_i)\ge L_0$. Thus $b_i\ge\beta$ and $a_i\le A_0$.
The quadratic penalty is at most $-L_0$, whence
\[
 |q-c|\le\sqrt{-2L_0/m},\quad
 \norm w\le\sqrt{-2L_0/\alpha},\quad b_i\le B_0.
\]
In particular $B_0>0$.

Let $f_i(z)=e^{a_i z}\mathbf1_{z\le b_i}$,
$\mu_i=\int f_i\,d\gamma_1=e^{\Psi(b_i)}<1$, and $h_i=f_i/\mu_i$.
They are centered in Gaussian first moment. For
\[
 (Dh)(u)=\int h((u+v)/\sqrt2)h((u-v)/\sqrt2)\,d\gamma_1(v),
\]
write $Q_0=\E_R\prod_i h_i(X_i)$ and $Q_1=\E_R\prod_i(Dh_i)(X_i)$.
Taking independent $X,X'\sim N(0,R)$ and setting
$U=(X+X')/\sqrt2$, $V=(X-X')/\sqrt2$, the vectors $U,V$ are independent
with covariance $R$. Conditional on $U=u$, the factors become
\begin{equation}\label{eq:intervalfactor}
 h_i((u_i+v_i)/\sqrt2)h_i((u_i-v_i)/\sqrt2)
 =\mu_i^{-2}e^{\sqrt2 a_i u_i}
   \mathbf1_{|v_i|\le\sqrt2 b_i-u_i}.
\end{equation}
On $\mathcal B=\{|U_i|\le\ell\ \forall i\}$ all radii belong to $[\ell,T]$.
The product of the weights in \eqref{eq:intervalfactor} is at least
$e^{-\sqrt2 mA_0\ell}=e^{-mA_0\beta}$. The symmetric rectangle inequality
gives $\Pp(\mathcal B)\ge b(\ell)^m$, even for singular $R$.
Lemma~\ref{lem:rectangle} therefore yields
\begin{equation}\label{eq:selfgap}
 Q_0^2-Q_1\ge K_{m,[u,v]}\norm{R-I_m}_F^2.
\end{equation}
Outside $\mathcal B$ the conditional difference remains nonnegative:
use the ordinary rectangle inequality when all radii are positive; if a
radius is negative both products vanish, and a zero radius also has zero
mass because every coordinate marginal is continuous.

Each $Dh_i$ is bounded and log-concave, has unit Gaussian integral and zero
Gaussian first moment. The integral and moment assertions follow by the same
orthogonal change of two independent scalar Gaussian variables; log-concavity
follows by integration of a log-concave function. Thus \eqref{eq:product} gives
$Q_1\ge1$. The Gaussian change of measure and \eqref{eq:tilt} give the exact identity
\[
 F_R(c)=e^{\mathcal J_*}Q_0.
\]
This identity is valid for singular $R$ as well: write $X=TZ$ with
$Z$ standard Gaussian in the rank of $R$ and tilt by
$T^{\mathsf T}a$. The shifted vector is $X+Ra$, and $b-Ra=c\one$;
no inverse covariance or full-dimensional density is used.
Consequently $F_R(c)\ge\Phi(c)^m\sqrt{1+K_{m,[u,v]}\norm{R-I_m}_F^2}$.
Subtract $\Phi(c)^m$ and use $\Phi(c)^m\ge\Phi(u)^m$.
Finally $\norm{R-I_m}_F^2\le m(m-1)$, since $|R_{ij}|\le1$.
Rationalizing the square root proves \eqref{eq:global}.
\end{proof}

\section{A centroid estimate and decoding stability}
For the rest of the global consequences use the interval $[0,1]$ and abbreviate
$\kappa=\kappa_{m,[0,1]}$. Define
\begin{equation}\label{eq:eta}
 M_0=\max_{y\in[H(0)/2,\,3H(1)/2]}(-\Psi''(y)),\quad
 \tau=\min\left\{\frac{H(0)}{2\alpha m},
                    \frac{r(1)}{1+M_0\alpha m}\right\},\quad
 \eta=2^{-m}\frac{\alpha r(1)\tau}{2}>0.
\end{equation}

\begin{lemma}[Centroid penalty]\label{lem:centroid}
For unit vectors $x_i$ with Gram matrix $G$ and $R$ as in
\eqref{eq:normalization}, set $s_X=\sum_i x_i$. For $0\le c\le1$,
\begin{equation}\label{eq:centroid}
 D_R(c)\ge\eta\norm{s_X}^2.
\end{equation}
\end{lemma}
\begin{proof}
In the variational functional use $q=H(c)$ and $w=\tau s_X$.
Since $\norm{s_X}\le m$, the scalar displacements
$z_i=\alpha\tau\langle x_i,s_X\rangle$ obey $|z_i|\le H(0)/2$.
Taylor's theorem on the interval in \eqref{eq:eta} gives
\[
 \Psi(H(c)+z_i)\ge\Psi(H(c))+r(c)z_i-M_0z_i^2/2.
\]
Now $\sum_i z_i=\alpha\tau\norm{s_X}^2$ and
$\sum_i z_i^2\le\alpha^2\tau^2m\norm{s_X}^2$.
Subtracting the value at $(H(c),0)$ gives
\[
 \mathcal J_* -m\log\Phi(c)
 \ge\alpha\left(r(c)\tau-\frac{\tau^2}{2}(1+M_0\alpha m)\right)\norm{s_X}^2
 \ge\frac{\alpha r(1)\tau}{2}\norm{s_X}^2.
\]
Use $F_R(c)\ge e^{\mathcal J_*}$, $e^z-1\ge z$, and $\Phi(c)^m\ge2^{-m}$.
\end{proof}

Let $\mu=\lambda/\sqrt\alpha$ and put
\begin{equation}\label{eq:codeconstants}
 W_{m,\lambda}=\frac{e^{-\mu^2/2}}m(e^\mu-1),\qquad
 c_g=W_{m,\lambda}\kappa\alpha^2,\quad
 c_s=W_{m,\lambda}\eta,\quad
 c_P=\left(\frac\alpha{c_g}+\frac1{mc_s}\right)^{-1}.
\end{equation}
These constants are positive for every fixed $m\ge2$ and $\lambda>0$.

\begin{theorem}[Quantitative recovery of a simplex]\label{thm:code}
Let $G$ be the Gram matrix of $m$ unit signals and
$\delta_\lambda(G)=\Pc(G_\Delta;\lambda)-\Pc(G;\lambda)$. Then
\begin{equation}\label{eq:code-stability}
 \delta_\lambda(G)\ge c_g\norm{G-G_\Delta}_F^2,\qquad
 \delta_\lambda(G)\ge c_s\norm{\textstyle\sum_i x_i}^2.
\end{equation}
If the ambient dimension is at least $m-1$, fix a regular simplex matrix
$V=[v_1\ \cdots\ v_m]$ in that space, and define
\[
 d_P(X,\Delta)^2=\min_{U\in O(d)}\sum_i\norm{x_i-Uv_i}^2.
\]
Then
\begin{equation}\label{eq:procrustes-stability}
 \delta_\lambda(G)\ge c_P d_P(X,\Delta)^2.
\end{equation}
In particular $\delta_\lambda(G)\le\varepsilon$ implies
$d_P(X,\Delta)\le\sqrt{\varepsilon/c_P}$.
\end{theorem}
\begin{proof}
Add an independent scalar $B\sim N(0,1/n)$ and set
$X_i'=\sqrt\alpha(\xi_i+B)$. Its covariance is $R$ and
\[
 \E e^{\mu\max_i X_i'}=e^{\lambda^2/(2n)}\E e^{\lambda\max_i\xi_i}.
\]
For the simplex the transformed coordinates are independent standard normals.
All the exponential moments are finite. Integrating their tail probabilities gives
\[
 \E e^{\mu\max_i Z_i}-\E e^{\mu\max_i X_i'}
 =\mu\int_{\R}e^{\mu c}D_R(c)\,dc.
\]
The integrand is nonnegative everywhere by the known comparison. Restrict it
to $[0,1]$, use Theorem~\ref{thm:global} or Lemma~\ref{lem:centroid}, and then
use \eqref{eq:pc}. The prefactor becomes
$m^{-1}e^{-\lambda^2/2-\lambda^2/(2n)}=m^{-1}e^{-\mu^2/2}$.
This proves \eqref{eq:code-stability}.

For completeness, the geometric conversion is quantitative linear algebra.
Let $X=[x_1\ \cdots\ x_m]$ and let $\lambda_1,\ldots,\lambda_n$ be the
eigenvalues of $PGP$ on $\one^\perp$, including zeros. The cross matrix $XV^{\mathsf T}$ has nonzero singular values
$\sqrt{\lambda_j/\alpha}$, with zeros added as necessary.
The orthogonal Procrustes formula follows by maximizing
$\tr(U^{\mathsf T}XV^{\mathsf T})$ over $O(d)$.
Alignment and the orthogonality of centered and constant columns give
\begin{align*}
 d_P(X,\Delta)^2
 &=\frac{\norm{s_X}^2}{m}
       +\sum_{j=1}^n(\sqrt{\lambda_j}-\alpha^{-1/2})^2\\
 &\le\frac{\norm{s_X}^2}{m}
       +\alpha\norm{PGP-G_\Delta}_F^2
 \le\frac{\norm{s_X}^2}{m}+\alpha\norm{G-G_\Delta}_F^2.
\end{align*}
The first inequality uses
$(\sqrt t-\sqrt a)^2\le(t-a)^2/a$ for $t\ge0,a>0$.
Applying the two separate deficit bounds proves \eqref{eq:procrustes-stability}.
Relabeling is already included: every permutation of regular-simplex vertices
extends to an orthogonal symmetry.
\end{proof}

\begin{corollary}[Insufficient dimension and unused energy]\label{cor:operational}
If the unit-signal Gram matrix has rank at most $d<m-1$, with $d\ge1$, then
\begin{equation}\label{eq:rank}
 \delta_\lambda(G)\ge c_gm^2\left(\frac1d-\frac1{m-1}\right).
\end{equation}
For arbitrary signals $\norm{x_i}\le1$, define the deficit using the actual
maximum-likelihood success probability (not the equal-energy formula
\eqref{eq:pc}). Then, regardless of whether the original channel permits an additional
use,
\begin{equation}\label{eq:slack}
 \Pc(G_\Delta;\lambda)-\Pc(\{x_i\};\lambda)
 \ge c_s\sum_i(1-\norm{x_i}^2).
\end{equation}
\end{corollary}
\begin{proof}
Using $\tr G=m$,
\[
 \norm{G-G_\Delta}_F^2
 =\tr(G^2)-\frac{m^2}{m-1}+\frac{2}{m-1}\one^{\mathsf T}G\one
 \ge\frac{m^2}{d}-\frac{m^2}{m-1}.
\]
Here $\tr(G^2)\ge(\tr G)^2/d$ and $\one^{\mathsf T}G\one\ge0$.
For the energy claim, use the following lift only as a mathematical
comparison (it imposes no additional-channel-use assumption on the code):
lift the codewords to
$\widetilde x_i=(x_i,q_i)$, $q_i=\sqrt{1-\norm{x_i}^2}\ge0$.
The augmented signals have unit norm, and their optimal success probability
is no smaller because the receiver may ignore the additional observation.
Apply the centroid inequality in \eqref{eq:code-stability} and note
$\norm{\sum_i\widetilde x_i}^2\ge(\sum_iq_i)^2\ge\sum_iq_i^2$.
\end{proof}

\section{Mean width of arbitrary inscribed simplices}
Set
\begin{equation}\label{eq:meanconstants}
 b_g=\kappa\alpha^{3/2},\qquad b_s=\eta/\sqrt\alpha.
\end{equation}
For any correlation matrix $G$ the same distribution-function integration,
without the exponential weight, proves
\begin{equation}\label{eq:mean-unit}
 \E\max_i\xi_i^\Delta-\E\max_i\xi_i
 \ge b_g\norm{G-G_\Delta}_F^2,\qquad
 \E\max_i\xi_i^\Delta-\E\max_i\xi_i
 \ge b_s\norm{\textstyle\sum_i x_i}^2.
\end{equation}
Indeed the expectation gap of the transformed maxima is
$\int_\R D_R(c)\,dc$, and the common added Gaussian has mean zero.

\begin{theorem}[Mean-width stability in a ball]\label{thm:width}
Let $y_1,\ldots,y_m\in B_2^d$ with $d\ge m-1$, let
$K=\operatorname{conv}\{y_i\}$, and let $\Delta$ be a unit regular simplex
in the same ambient space. Use normalized spherical mean width, so that
$w(K)=2\E\max_i\langle g,y_i\rangle/\E\norm g$ for $g\sim N(0,I_d)$.
Then
\begin{equation}\label{eq:width-stability}
 w(\Delta)-w(K)
 \ge \frac{2}{\E\norm g}
          \left(\frac\alpha{b_g}+\frac1{b_s}\right)^{-1}
          d_P([y_i],\Delta)^2.
\end{equation}
Also the unnormalized expected-maximum deficit is at least
$b_s\sum_i(1-\norm{y_i}^2)$.
\end{theorem}
\begin{proof}
Put $d_i=1-\norm{y_i}^2$ and lift to unit vectors
$\widehat y_i=(y_i,\sqrt{d_i}e_i)$ using mutually orthogonal new coordinates.
Conditional Jensen implies that adding these independent mean-zero noises
cannot decrease the expected maximum. Write $G=(\langle y_i,y_j\rangle)$,
$\widehat G=G+\diag(d_i)$,
$T_g=\norm{\widehat G-G_\Delta}_F^2$, and
$T_s=\norm{\sum_i y_i}^2+\sum_i d_i$.
The original expected-maximum deficit $\delta_w$ is at least $b_gT_g$ and
at least $b_sT_s$ by \eqref{eq:mean-unit}. The Procrustes argument in the
previous proof applies to vectors of any lengths, and hence
\begin{align*}
 d_P([y_i],\Delta)^2
 &\le\alpha\norm{G-G_\Delta}_F^2+\norm{\textstyle\sum_i y_i}^2/m\\
 &=\alpha\left(T_g+\sum_i d_i^2\right)+\norm{\textstyle\sum_i y_i}^2/m
 \le\alpha T_g+T_s,
\end{align*}
since $0\le d_i\le1$. Thus
$d_P^2\le(\alpha/b_g+1/b_s)\delta_w$.
The polar decomposition of a standard Gaussian gives \eqref{eq:width-stability}.
\end{proof}

\section{Sharp local coefficients and optimal exponents}
The explicit global constants are deliberately conservative. The exact local
coefficient is much simpler and identifies a distinction between $m=3$ and
$m\ge4$.

\begin{theorem}[Sharp local common-threshold coefficient]\label{thm:local}
Fix a finite $c$, and write $p=\Phi(c)$, $q=r(c)$, $h=c+q$, and
$k=\phi(c)^2p^{m-2}$. For $m\ge3$,
\begin{equation}\label{eq:sharp}
 \liminf_{\substack{R\to I_m,\ R\in\mathcal A_m\\R\ne I_m}}
   \frac{D_R(c)}{\norm{R-I_m}_F^2}
 =\begin{cases}
    k(1+ch)/4,&m=3,\\
    kh^2/4,&m\ge4.
   \end{cases}
\end{equation}
Both constants are positive. For $m=2$ the sharp behavior is linear:
$D_R(c)/\norm{R-I_2}_F\to\phi(c)^2/\sqrt2$.
\end{theorem}
\begin{proof}
Put $A=R-I_m$, $S=\sum_{i<j}A_{ij}$,
$Q=\sum_{i<j}A_{ij}^2=\norm A_F^2/2$, and $r_A=A\one$.
The Gaussian density is smooth in $A$ near zero and its first three derivatives
are bounded by an integrable polynomial times a Gaussian. Differentiation
under the integral over the fixed lower orthant is therefore justified.
For off-diagonal covariance entries, the first derivative at independence is
$k$. The second-order Taylor contributions for a repeated edge, two edges
sharing one endpoint, and two disjoint edges are respectively
\[
 \tfrac12 c^2\phi(c)^2p^{m-2}A_{ij}^2,\quad
 -c\phi(c)^3p^{m-3}A_{ij}A_{ik},\quad
 \phi(c)^4p^{m-4}A_{ij}A_{k\ell}.
\]
Each unordered pair of distinct edges is counted once. The sums of products
over adjacent and disjoint pairs are
$\norm{r_A}^2/2-Q$ and $(S^2+Q-\norm{r_A}^2)/2$, respectively.
Consequently
\begin{equation}\label{eq:taylor}
 \frac{D_{I_m+A}(c)}k
 =S+\frac{h^2}{2}Q-\frac{qh}{2}\norm{r_A}^2
       +\frac{q^2}{2}S^2+O(\norm A_F^3).
\end{equation}
For $m=3$ there are no disjoint pairs; the displayed combined formula still
holds because the corresponding combinatorial sum is zero.

Feasibility means $L=P+A\succeq0$. Since $L\one=r_A$, the spectral inequality
$L^2\preceq\lambda_{\max}(L)L$ gives
\begin{equation}\label{eq:feasible}
 S\ge0,\qquad \norm{r_A}^2\le2(1+\norm A_{\mathrm{op}})S.
\end{equation}
Also $\sigma=1-qh>0$ is a truncated-normal variance.
Equations \eqref{eq:taylor}--\eqref{eq:feasible} imply, for sufficiently small
feasible $A$,
\[
 D_{I_m+A}(c)\ge k\left(\frac\sigma2 S+\frac{h^2}{4}Q\right).
\]
Thus sequences with $S/Q\to\infty$ have an infinite deficit quotient.
On sequences with $S=O(Q)$, \eqref{eq:feasible} yields
$S\ge\norm{r_A}^2/2+o(Q)$ and hence
\begin{equation}\label{eq:local-lower}
 D_{I_m+A}(c)/k\ge\tfrac12 h^2Q
                       +\tfrac12\sigma\norm{r_A}^2+o(Q).
\end{equation}
For $m\ge4$ this gives the lower coefficient $kh^2/4$.
It is attained by $A=tD$, where $D$ is symmetric with
$D_{12}=D_{34}=1$, $D_{23}=D_{14}=-1$, and all other entries zero.
Then $D\one=0$, so $P+tD\succeq0$ for sufficiently small $|t|$.

For $m=3$ there is the identity $\norm{r_A}^2=Q+S^2$.
Equation \eqref{eq:local-lower} gives the lower coefficient
$k(h^2+\sigma)/4=k(1+ch)/4$.
To attain it, take a regular simplex $v_i\in\R^{m-1}$, a unit vector
$u\perp v_1$, and replace only $v_1$ by $v_1(t)=\cos(t)v_1+\sin(t)u$.
For the resulting $A=\alpha(G(t)-G_\Delta)$ one computes exactly
\[
 S=\alpha(1-\cos t),\quad
 Q=\frac{\alpha^2}{m-1}\big((1-\cos t)^2+m\sin^2t\big),\quad
 \norm{r_A}^2=2S.
\]
Thus $Q=\alpha t^2+O(t^4)$ and the desired quotient follows from
\eqref{eq:taylor}. Positivity follows from $h^2+\sigma>0$.
For $m=2$, feasibility is $R_{12}=\rho\ge0$ and the first derivative at
$\rho=0$ is $\phi(c)^2$, whereas $\norm{R-I_2}_F=\sqrt2\rho$.
\end{proof}

\begin{corollary}[Optimal powers]\label{cor:optimal}
For $m\ge3$, the power two of the Frobenius distance in the common-threshold
bound and in decoding Gram stability cannot be replaced by any smaller
positive power with a uniform positive constant near the simplex.
For $m\ge3$, the power two of the coordinate distance in \eqref{eq:procrustes-stability}
is likewise optimal, including when signals are restricted to dimension $m-1$.
\end{corollary}
\begin{proof}
Use the one-vertex rotation from the proof of Theorem~\ref{thm:local}.
It gives $\norm{G(t)-G_\Delta}_F=\Theta(|t|)$ and $S=O(t^2)$.
The normalized covariance $R(t)$ remains near $I_m$, where expectations of
$e^{\mu\max_i X_i}$ are twice differentiable by Gaussian-density domination.
Permutation symmetry makes their first derivative a constant times the sum
of the off-diagonal perturbations. That sum has zero first derivative along
this path. The decoding deficit is therefore $O(t^2)$, and the lower bound
in Theorem~\ref{thm:code} makes it $\Theta(t^2)$.
The same argument applies to the fixed-threshold deficit.
Moreover $\norm{G-G_\Delta}_F\le2\sqrt m\,d_P(X,\Delta)$ for unit columns,
while the displayed rotation itself gives $d_P=O(|t|)$. Hence $d_P=\Theta(|t|)$.
No estimate with a power below two can hold on this path.
\end{proof}

\section{Second variation in the configuration metric}\label{sec:spectrum}

The Frobenius normalization in Theorem~\ref{thm:local} is different from the
metric of signal configurations.  We record the distinction explicitly.
The calculation is a reformulation of local second variation, not a
claim of priority for local optimality or the existence of a
second-variation calculation.

We first complete the minimum-direction statement in the Gram metric.
For a feasible sequence $A=R-I_m\to0$, let $\zeta\ge0$ be the Schur slack
defined in Lemma~\ref{rem:quotient}. For $m\ge4$ the sharp quotient in
Theorem~\ref{thm:local} is attained exactly when
$\|A\one\|^2+\zeta=o(\|A\|_F^2)$. Thus its nonzero limiting directions
are all symmetric zero-diagonal, zero-row-sum matrices, with no excess
second-order Schur slack. For $m=3$ the corresponding condition is
$\zeta=o(\|A\|_F^2)$: all nonzero critical tangent directions can attain
the minimum. These conditions follow from the positive coefficients of
$\|A\one\|^2$ and $\zeta$ in the expansions below, which specify the
second-order feasibility condition as well as the first-order direction.

Fix a unit regular simplex $V=[v_1\ \cdots\ v_m]$ in $\R^d$, $d\ge n=m-1$,
and set $E=\operatorname{span}\{v_i\}$ and $r=d-n$.  The configuration space
is $(S^{d-1})^m$, with squared tangent norm $\|U\|_F^2=\sum_i\|u_i\|^2$.
At $V$ the rotation directions are $U=\Omega V$, $\Omega^{\mathsf T}=-\Omega$.
Their orthogonal complement will be called the horizontal tangent space.
Write
\[
 \mathcal H=\{B=B^{\mathsf T}:\diag B=0,\ \one^{\mathsf T}B\one=0\}.
\]
For $m\ge3$ this space has the orthogonal decomposition
\[
 \mathcal H=\mathcal H_0\oplus\mathcal H_1,\qquad
 \mathcal H_0=\{B\in\mathcal H:B\one=0\},
\]
where $\mathcal H_1$ consists of the matrices
$B_{ij}=a_i+a_j$ for $i\ne j$, $B_{ii}=0$, and $a\perp\one$.
The dimensions are $m(m-3)/2$ and $m-1$, respectively. For $m=3$, $\mathcal H_0=
\{0\}$.

\begin{theorem}[Horizontal second variation and its minimum]\label{thm:spectrum}
Fix a finite $c$, and retain $k,h,q$ from Theorem~\ref{thm:local}; put
$\sigma=1-qh>0$.  For $m\ge3$ the horizontal tangent space is the orthogonal
sum of the lifts of $\mathcal H_0$, the lifts of $\mathcal H_1$, and the
normal-centroid space $\{z\one^{\mathsf T}:z\in E^\perp\}$.
For a twice differentiable sphere path with horizontal initial velocity
$U=U_0+U_1+U_\perp$,
\begin{equation}\label{eq:spectrum-cdf}
 D_{R(t)}(c)=t^2\left(
 k\alpha h^2\|U_0\|_F^2+
 \frac{k}{2}\bigl[h^2+(m-2)\sigma\bigr]\|U_1\|_F^2+
 \frac{k(m-1)}2\|U_\perp\|_F^2\right)+o(t^2).
\end{equation}
Thus the Hessian eigenvalues are twice the three displayed coefficients,
with multiplicities $m(m-3)/2$, $m-1$, and $r$, omitting zero-dimensional
spaces.  On the full tangent space there is also the zero eigenvalue, of
multiplicity $n(n-1)/2+nr$, precisely the rotation directions.

The exact local coefficient of $D_R(c)$ relative to $d_P(X,\Delta)^2$ is
the minimum of the coefficients in \eqref{eq:spectrum-cdf} whose spaces
are present.  All minimum directions are precisely the direct sum of the
spaces with that minimum coefficient, modulo rotations and the residual
simplex symmetries.  In particular:
\begin{itemize}
\item If $m\ge4$ and $d=n$, the minimizing space is $\mathcal H_0$ when
$h^2<m\sigma$, $\mathcal H_1$ when $h^2>m\sigma$, and their direct sum
when equality holds.
\item If $m\ge4$ and $d\ge m$, the minimum is
$k\alpha\min\{h^2,m/2\}$.  The minimizing space is $\mathcal H_0$ when
$h^2<m/2$, the normal-centroid space when $h^2>m/2$, and their direct sum
when equality holds.  The $\mathcal H_1$ coefficient is strictly larger
than this minimum for every finite $c$.
\item If $m=3$ and $d=2$, all nonzero horizontal directions minimize,
with coefficient $k(1+ch)/2$.  If $m=3$ and $d\ge3$, compare this
coefficient to $k$: the two minimizing spaces are selected by $ch<1$,
$ch>1$, or their direct sum at $ch=1$.
\end{itemize}
For $m=2$ and $d\ge2$, only the normal-centroid space remains, with
coefficient $\phi(c)^2/2$ and horizontal multiplicity $d-1$.  For $m=2$,
$d=1$, the configuration space is locally discrete and there is no
nonconstant local variation.
\end{theorem}

\begin{proof}
Let $X(t)$ be a sphere path with $X(0)=V$, $X'(0)=U$, and
$A(t)=\alpha(X(t)^{\mathsf T}X(t)-G_\Delta)$.  Its initial Gram velocity is
\[
 B=A'(0)=\alpha(V^{\mathsf T}U+U^{\mathsf T}V)\in\mathcal H.
\]
Since $S(t)=\alpha\|X(t)\one\|^2/2$ exactly,
\[
 S(t)=\frac\alpha2t^2\|U\one\|^2+o(t^2),\qquad
 B\one=\alpha V^{\mathsf T}U\one.
\]
Using $VV^{\mathsf T}=\alpha^{-1}I_E$ and separating the normal projection
$U_N$ of $U$, equation~\eqref{eq:taylor} therefore gives
\begin{equation}\label{eq:intrinsic-quadratic}
 \frac{D_{R(t)}(c)}{kt^2}
 =\frac{h^2}{4}\|B\|_F^2+
   \frac\sigma2\|B\one\|^2+
   \frac\alpha2\|U_N\one\|^2+o(1).
\end{equation}
In particular the first variation vanishes and the second variation does
not depend on acceleration.

Every $B\in\mathcal H$ has the unique horizontal in-plane lift
\[
 U_B=\frac12VB+\frac1{2m}V(B\one)\one^{\mathsf T}.
\]
Direct multiplication proves its tangency, horizontality, and prescribed
Gram velocity.  Its squared norm is
\begin{equation}\label{eq:horizontal-metric}
 \|U_B\|_F^2
 =\frac1{4\alpha}\|B\|_F^2+
   \frac1{2m\alpha}\|B\one\|^2.
\end{equation}
The remaining horizontal vectors are exactly $z\one^{\mathsf T}$ with
$z\in E^\perp$: all normal matrices with zero column sum are rotation
directions.  The two in-plane spaces are orthogonal both in Frobenius
Gram norm and in \eqref{eq:horizontal-metric}.  On $\mathcal H_0$ the
metric is $\|B\|_F^2/(4\alpha)$.  On $\mathcal H_1$,
\[
 \|B\one\|^2=\frac{m-2}{2}\|B\|_F^2,\qquad
 \|U_B\|_F^2=\frac12\|B\|_F^2.
\]
For $U_\perp=z\one^{\mathsf T}$ one has
$\|U_\perp\one\|^2=m\|U_\perp\|_F^2$.
Substitution proves the spectrum.  Every listed direction is attainable:
use the independent spherical geodesics
\[
 x_i(t)=\cos(t\|u_i\|)v_i+
       \frac{\sin(t\|u_i\|)}{\|u_i\|}u_i,
\]
with the continuous interpretation when $u_i=0$.  Furthermore
$d_P(X(t),\Delta)^2=t^2\|U\|_F^2+o(t^2)$ for horizontal $U$.
The path-free argument below justifies the corresponding local infimum
over arbitrary sequences, rather than only over differentiable paths.

To verify the stated ordering when $m\ge4$ and the normal space is
present, use $h'=\sigma$ and $\sigma'=q(h^2-\sigma)$.  Put
$f=\sigma-h^2/m$.  At a zero of $f$,
\[
 f'=\frac{h}{m^2}\bigl(m(m-1)-(m+1)h^2\bigr)>0
 \quad\hbox{if }h^2\le m/2.
\]
The elementary Gaussian-tail asymptotics give $\sigma/h^2\to1$ as
$c\to-\infty$.  Hence a first downward crossing is impossible and
$\sigma>h^2/m$ whenever $h^2\le m/2$.  At $h^2=m/2$ this gives
$\sigma>1/2$.  Subsequently $\sigma$ cannot cross down through $1/2$,
because $\sigma'=q(h^2-1/2)>0$ on that boundary.  These two strict
inequalities imply that the $\mathcal H_1$ coefficient exceeds the
smaller of the other two coefficients.  The other comparisons follow
by subtraction; for $m=3$, $h^2+\sigma=1+ch$.

For $m=2$, write $A_{12}=\rho\ge0$.  Then
$D_R(c)=\phi(c)^2\rho+O(\rho^2)$ and the exact Procrustes formula is
$d_P^2=4(1-\sqrt{1-\rho})=2\rho+O(\rho^2)$.
\end{proof}

\begin{lemma}[Feasible-sequence expansion]\label{rem:quotient}
Let $m\ge3$ and let $X$ be a unit-column configuration approaching the
simplex orbit in $\R^d$, with $d\ge n$. Put $A=R-I_m$, $e=\one/\sqrt m$
and $C=PAP|_{e^\perp}$. Define
\[
 \zeta=e^{\mathsf T}Ae-(PAe)^{\mathsf T}(I+C)^{-1}(PAe),\qquad
 B=A-\frac{2S}{m(m-1)}(J-I_m),\qquad b=B\one.
\]
For $i\ne j$, set $(B_1)_{ij}=(b_i+b_j)/(m-2)$, set $(B_1)_{ii}=0$
and put $B_0=B-B_1$. Then $\zeta\ge0$, $B_j\in\mathcal H_j$ and
\[
 E_0=\frac{\|B_0\|_F^2}{4\alpha},\qquad
 E_1=\frac{\|B_1\|_F^2}{2},\qquad E_\perp=\frac\zeta\alpha
\]
satisfy
\begin{align}
 d_P^2&=E_0+E_1+E_\perp+o(d_P^2),\label{eq:sequence-metric}\\
 D_R(c)&=k\alpha h^2E_0+
 \frac{k}{2}[h^2+(m-2)\sigma]E_1+
 \frac{k(m-1)}2E_\perp+o(d_P^2).\label{eq:sequence-deficit}
\end{align}
Here $\zeta=0$ if and only if $\rank G=n$. In particular $E_\perp=0$
when $d=n$. The remainders apply to arbitrary feasible sequences.
\end{lemma}
\begin{proof}
In the orthogonal splitting $\R e\oplus e^\perp$, the matrix $P+A$ has blocks
\[
 P+A=\begin{pmatrix}a&v^{\mathsf T}\\v&I+C\end{pmatrix},
 \qquad a=2S/m,\quad v=PAe.
\]
The lower block is positive definite near $A=0$. The Schur complement
therefore gives positivity exactly when $\zeta\ge0$ and rank $n$ exactly
when $\zeta=0$. Since $\|v\|^2=\|A\one\|^2/m-4S^2/m^2$, expansion of
$(I+C)^{-1}$ gives
\begin{equation}\label{eq:schur-sum}
 S=\frac12\|A\one\|^2+\frac m2\zeta+
 O(\|A\|_F^3+S^2).
\end{equation}
The exact Procrustes formula from Theorem~\ref{thm:code} is
\[
 d_P^2=2m-\frac2\alpha\tr\sqrt{I+C}.
\]
Since $\tr C=-2S/m$, its matrix square-root expansion gives
\[
 d_P^2=\frac{2S}{m\alpha}+\frac{\|C\|_F^2}{4\alpha}
       +O(\|A\|_F^3).
\]
The block decomposition also gives
\[
 \|C\|_F^2=\|A\|_F^2-\frac2m\|A\one\|^2+
                   \frac{4S^2}{m^2}.
\]
Substitution of \eqref{eq:schur-sum} yields
\begin{equation}\label{eq:schur-metric}
 d_P^2=\frac\zeta\alpha+\frac{\|A\|_F^2}{4\alpha}
       +\frac{\|A\one\|^2}{2m\alpha}
       +O(\|A\|_F^3+S^2).
\end{equation}
To control this remainder in the configuration metric, align $V$ to a
minimizer of $d_P$. Then
$S=\alpha\|(X-V)\one\|^2/2\le\alpha m d_P^2/2$ and
$\|A\|_F\le2\alpha\sqrt m\,d_P$.
Thus the remainder is $o(d_P^2)$ uniformly for fixed $m$.
Moreover $A-B=O(S)$, so replacing $A$ by $B$ in the quadratic terms
changes them by $o(d_P^2)$.
Orthogonality of $B_0,B_1$ and
$\|B_1\one\|^2=(m-2)\|B_1\|_F^2/2$ now prove
\eqref{eq:sequence-metric}.
Substituting \eqref{eq:schur-sum} into \eqref{eq:taylor} gives
\[
 D_R(c)/k=\frac{h^2}{4}\|A\|_F^2+
       \frac\sigma2\|A\one\|^2+\frac m2\zeta+o(d_P^2).
\]
The same orthogonal decomposition proves \eqref{eq:sequence-deficit}.
\end{proof}

Because the three energies are nonnegative, a feasible sequence attains the
minimum quotient exactly when its higher-coefficient energies sum to
$o(d_P^2)$. This proves the arbitrary-sequence claim in
Theorem~\ref{thm:spectrum}. The rotation-and-permutation quotient itself
need not be smooth. Its tangent cone is
$(\mathcal H\oplus E^\perp)/(S_m\times O(r))$ with the horizontal metric.
For $r>0$, the normal-centroid component becomes a radial half-line in the
quotient, while its horizontal multiplicity remains $r$.

The distinction between Gram and configuration metrics is explicit when
$d\ge m$. For a unit normal $z\perp E$, set
$x_i(t)=\sqrt{1-t^2}\,v_i+tz$. Then
\[
 A=t^2(J-I_m),\qquad
 d_P^2=2m(1-\sqrt{1-t^2})=mt^2+O(t^4),\qquad
 \Big\|\sum_i x_i(t)\Big\|^2=m^2t^2.
\]
The squared Gram defect has order $t^4$, whereas the squared configuration
distance has order $t^2$. The separate centroid term controls this transverse path.

\begin{corollary}[Decoding and mean-width second variations]\label{cor:spectrum-integrated}
Let $\mu=\lambda/\sqrt\alpha$ and define
\[
 \omega_\lambda(c)=\frac{\mu e^{-\mu^2/2}}m
 e^{\mu c}\phi(c)^2\Phi(c)^{m-2},\qquad
 A_\lambda=\int_\R\omega_\lambda h^2,\quad
 B_\lambda=\int_\R\omega_\lambda\sigma,\quad
 C_\lambda=\int_\R\omega_\lambda.
\]
All three constants are finite and positive.  The quadratic coefficients
of the decoding deficit on the three horizontal spaces are exactly
\[
 \Lambda_0=\alpha A_\lambda,\qquad
 \Lambda_1=\frac{A_\lambda+(m-2)B_\lambda}{2},\qquad
 \Lambda_\perp=\frac{m-1}{2}C_\lambda.
\]
The local coefficient relative to $d_P^2$ is their minimum over spaces
present, and every minimum direction belongs to the direct sum of the
spaces attaining that minimum.  Theorems~\ref{thm:ordered-switches} and~\ref{thm:triangular-switch}
below determine the ordering of these decoding coefficients at every positive
channel parameter.

Relative to the Gram Frobenius norm, the exact decoding coefficient is
\[
 \liminf_{G\to G_\Delta,\ G\ne G_\Delta}
 \frac{\delta_\lambda(G)}{\|G-G_\Delta\|_F^2}
 =\begin{cases}
 \alpha^2(A_\lambda+B_\lambda)/4,&m=3,\\
 \alpha^2 A_\lambda/4,&m\ge4.
 \end{cases}
\]
These coefficients are attainable in dimension $n$.  For $m=2$,
\[
 \frac{\delta_\lambda(G)}{\|G-G_\Delta\|_F}
 \longrightarrow\frac{\lambda\phi(\lambda)}{4\sqrt2},
 \qquad
 \frac{\delta_\lambda(G)}{d_P^2}
 \longrightarrow\frac{\lambda\phi(\lambda)}4
\]
along nonconstant configurations in dimension at least two.

For $m\ge3$, the same assertions for the unnormalized expected-maximum deficit hold
after replacing $\omega_\lambda(c)$ by
$\omega_w(c)=\phi(c)^2\Phi(c)^{m-2}/\sqrt\alpha$.
For normalized spherical mean width, multiply all coefficients by
$2/\E\|g\|$, $g\sim N(0,I_d)$.
For $m=2$ in dimension at least two, the expected-maximum deficit has
local coefficients $\phi(0)/(2\sqrt2)$ relative to
$\|G-G_\Delta\|_F$ and $\phi(0)/2$ relative to $d_P^2$.
\end{corollary}
\begin{proof}
Use the integral identities already established in the manuscript.
Near $R=I_m$, covariance-density derivatives through order three are
bounded by a polynomial times a fixed nondegenerate Gaussian density.
Their products with $e^{\mu\max_i x_i}$, or with $\max_i x_i$, remain
integrable.  Thus differentiating the expectations twice is justified,
and their Taylor coefficients equal the displayed integrals.  Equivalently,
differentiate the distribution-function identities with a common
integrable domination: for negative $c$ bound the fixed orthant integral
by a Gaussian tail, and for positive $c$ use that each derivative density
has total integral zero and integrate over the complementary tail.
This also proves convergence of the coefficient integrals.

Integrating the path-free expansion in Lemma~\ref{rem:quotient}, or
applying this density Taylor expansion directly, proves the coordinate
coefficients and their minimum criterion.  For the Gram coefficients,
the expansion is
\[
 \delta_\lambda(G)
 =\frac{A_\lambda}{4}\|A\|_F^2+
   \frac{B_\lambda}{2}\|A\one\|^2+
   \frac{mC_\lambda}{2}\zeta+o(\|A\|_F^2)
\]
on sequences with $S=O(\|A\|_F^2)$.  Sequences for which
$S/\|A\|_F^2\to\infty$ have divergent deficit quotient, by the positive
centroid bound already proved.  For $m\ge4$ the lower coefficient is
attained by a nonzero row-zero Gram direction with zero Schur slack.
For $m=3$ use
$\|A\one\|^2=\|A\|_F^2/2+S^2$ and a rank-two feasible path.
Finally $\|A\|_F=\alpha\|G-G_\Delta\|_F$.
For $m=2$ the exact success probability is
$\Phi(\lambda\sqrt{1-\rho})$, $\rho=A_{12}$; the stated constants follow
by differentiation and the Procrustes formula above.
\end{proof}

\section{Unique channel transitions and exact small-shell minimizers}
\label{sec:new-results}

Corollary~\ref{cor:spectrum-integrated} expresses the directional coefficients through
three explicit integrals. We now determine the ordering of these coefficients for every positive
channel parameter and, when the normal coefficient is strictly least,
identify the exact optimizer on every sufficiently small Procrustes shell.
The radii obtained here are uniform in the ambient dimension, but are not
evaluated numerically.

\subsection{A tilted Gaussian family}

For $m\ge3$ and $\mu\in\R$, set
\begin{align}
 w_{m,\mu}(c)&=e^{\mu c}\phi(c)^2\Phi(c)^{m-2},&
 Z_m(\mu)&=\int_\R w_{m,\mu}(c)\,dc,\label{new:tilt}\\
 a_m(\mu)&=Z_m(\mu)^{-1}\int_\R h(c)^2w_{m,\mu}(c)\,dc,&
 b_m(\mu)&=Z_m(\mu)^{-1}\int_\R \sigma(c)w_{m,\mu}(c)\,dc,\notag\\
 H_m(\mu)&=Z_m(\mu)^{-1}\int_\R h(c)w_{m,\mu}(c)\,dc.\notag
\end{align}
Here $q=\phi/\Phi$, $h(c)=c+q(c)$ and $\sigma=1-qh$, as in
Theorem~\ref{thm:spectrum}. The normalizer is positive and finite.
At $\mu=\lambda/\sqrt\alpha$,
\[
 A_\lambda=C_\lambda a_m(\mu),\qquad
 B_\lambda=C_\lambda b_m(\mu),\qquad
 C_\lambda=\frac{\mu e^{-\mu^2/2}}m Z_m(\mu)>0.
\]
Consequently the normalized decoding curvatures are
\begin{equation}\label{new:normalized-curvatures}
 L_0=\frac{m-1}{m}a_m(\mu),\qquad
 L_1=\frac{a_m(\mu)+(m-2)b_m(\mu)}2,\qquad
 L_\perp=\frac{m-1}{2}.
\end{equation}
The endpoint $\mu=0$ below is an endpoint of the normalized moment family;
it is not division by $C_\lambda$ at a zero-amplitude channel.

\begin{theorem}[Unique ordered channel crossings]\label{thm:ordered-switches}
Let $m\ge4$. There are unique positive parameters
\[
 0<\mu_{1\perp}(m)<\mu_{0\perp}(m)<\mu_{01}(m)
\]
satisfying, respectively,
\[
 a_m+(m-2)b_m=m-1,\qquad a_m=m/2,\qquad a_m=mb_m.
\]
Each equation has no other root in $[0,\infty)$. For $\lambda>0$,
with $\mu=\lambda/\sqrt\alpha$, the four open intervals have the orderings
\begin{align*}
 0<\mu<\mu_{1\perp}&:\quad \Lambda_0<\Lambda_1<\Lambda_\perp,\\
 \mu_{1\perp}<\mu<\mu_{0\perp}&:\quad \Lambda_0<\Lambda_\perp<\Lambda_1,\\
 \mu_{0\perp}<\mu<\mu_{01}&:\quad \Lambda_\perp<\Lambda_0<\Lambda_1,\\
 \mu>\mu_{01}&:\quad \Lambda_\perp<\Lambda_1<\Lambda_0.
\end{align*}
At the three endpoints the indicated pair agrees, and the remaining
inequality is strict. In particular
$\min\{\Lambda_0,\Lambda_\perp\}<\Lambda_1$ for every $\lambda>0$.
The thresholds in $\lambda$ are $\sqrt\alpha$ times these thresholds.
\end{theorem}

\begin{proof}
Appendix~\ref{ap:transitions} proves continuity and strict increase of
$a_m$, $b_m$, $H_m$ and $a_m/b_m$, as well as
\[
 \frac12<b_m(\mu)<1\quad(m\ge4,\ \mu\ge0),\qquad
 2a_m+(m-1)b_m=m+\mu H_m.
\]
At zero the identity gives $a_m<m/2$ and, using $b_m<1$,
$a_m+(m-2)b_m<m-1$. Using $b_m>1/2$ also gives $a_m<mb_m$.
Moreover $2a_m(\mu)>1+\mu H_m(0)$, so $a_m$ tends to infinity.
The three continuous strictly increasing functions
$a_m+(m-2)b_m$, $a_m$, and $a_m/b_m$ therefore cross their targets once.
At $a_m=m/2$, the bound $b_m>1/2$ puts the first crossing to the left
and the third to the right. Finally
$L_1-L_0=(m-2)(mb_m-a_m)/(2m)$ gives the displayed sign chart.
The analytic ingredients, including the strict benchmark at $m=4$,
are proved in full in Appendix~\ref{ap:transitions}.
\end{proof}

\begin{theorem}[The triangular channel crossing]\label{thm:triangular-switch}
For $m=3$ there is a unique $\mu_*(3)>0$ such that
$\mu_*H_3(\mu_*)=1$, equivalently $a_3(\mu_*)+b_3(\mu_*)=2$.
For $0<\mu<\mu_*$ one has $\Lambda_1<\Lambda_\perp$; for
$\mu>\mu_*$ the reverse inequality holds. They agree only at $\mu_*$.
The space $\mathcal H_0$ is absent.
\end{theorem}
\begin{proof}
The balance identity gives
$(a_3+b_3)/2=(3+\mu H_3)/4$.
On $[0,\infty)$, the function $\mu H_3(\mu)$ is continuous and strictly
increasing from zero to infinity, since $H_3$ is positive and increasing.
This proves the assertion. Its integral ingredients are included in
Appendix~\ref{ap:transitions}.
\end{proof}

\subsection{The exact strict-normal shell optimum}

For $m\ge3$ and $\lambda>0$, impose the strict-normal condition
\begin{equation}\label{new:normal-gap}
 \begin{cases}
  \mu>\mu_{0\perp}(m),&m\ge4,\\
  \mu>\mu_*(3),&m=3,
 \end{cases}
 \qquad \mu=\lambda/\sqrt\alpha.
\end{equation}
Equivalently $a_m(\mu)>m/2$ for $m\ge4$, and
$a_3(\mu)+b_3(\mu)>2$ for $m=3$. These are precisely the conditions that
the normal coefficient be strictly smaller than every in-plane coefficient
whose space is present.

\begin{theorem}[Exact optimizer on small Procrustes shells]\label{thm:normal-shell}
Assume \eqref{new:normal-gap}. There is $\rho_0=\rho_0(m,\lambda)>0$
such that, for every $d\ge m$, every unit regular simplex
$V=[v_1\ \cdots\ v_m]$ in $\R^d$, and every $0<\rho<\rho_0$, the minimum
of $\delta_\lambda(X)=\Pc(G_\Delta;\lambda)-\Pc(G_X;\lambda)$ over all
unit configurations with $d_P(X,V)=\rho$ is attained precisely by the
orthogonal images of
\begin{equation}\label{new:normal-minimizer}
 N_i=\kappa v_i+\sqrt{1-\kappa^2}\,z,\qquad
 \kappa=1-\frac{\rho^2}{2m},\qquad
 \|z\|=1,\quad z\perp\operatorname{span}\{v_i\}.
\end{equation}
The normal $z$ may be chosen once for all these radii. Different choices
of $z$ give the same orthogonal orbit. The minimal deficit is exactly
\begin{equation}\label{new:exact-deficit}
 \delta_{\min}(\rho)=
 \Pc(G_\Delta;\lambda)-
 \Pc\left(G_\Delta;\lambda\left(1-\frac{\rho^2}{2m}\right)\right).
\end{equation}
\end{theorem}

\begin{proof}
Appendix~\ref{ap:normal} supplies the complete geometric argument.
The centered square root
$T=(\alpha PG_XP)^{1/2}$ satisfies, on a uniform neighborhood,
$d_P(X,V)^2=2m-2\tr(T)/\alpha$.
On the shell $\tau=\rho^2$, write $T=\kappa P+D$, where
$D$ is symmetric, centered and trace-free.
The covariance reconstructed from $(\tau,D)$ defines a smooth
extension of the decoding deficit. Permutation symmetry makes $D=0$
stationary for every small $\tau$.
Its Hessian at $(0,0)$ is positive definite under \eqref{new:normal-gap};
the precise coefficients and the square-root coordinate normalization
are derived in Appendix~\ref{ap:normal}.
Continuity then gives a strict quadratic lower bound near this ray.
Every sufficiently small actual shell lies in that neighborhood.
Thus $D=0$ is its unique minimizing covariance.
It is realized by \eqref{new:normal-minimizer}, with exact distance $\rho$,
and equality of Gram matrices is equivalent to an orthogonal change
of realization. A common Gaussian shift in the normal direction gives
\eqref{new:exact-deficit}. The neighborhood estimates are made in the
fixed $m\times m$ Gram space, so $\rho_0$ is independent of $d$ and $V$.
\end{proof}

\begin{theorem}[Quartic expansion with a sixth-order remainder]
\label{thm:normal-taylor}
Let $p(\lambda)=\Pc(G_\Delta;\lambda)$. This function is real analytic.
For fixed $m\ge3$ and $\lambda>0$ satisfying \eqref{new:normal-gap},
there are $C=C(m,\lambda)>0$ and $r=r(m,\lambda)>0$ such that
Theorem~\ref{thm:normal-shell} holds for $0<\rho<r$ and
\begin{align}
 \delta_{\min}(\rho)
 &=\frac{\lambda p'(\lambda)}{2m}\rho^2
   -\frac{\lambda^2p''(\lambda)}{8m^2}\rho^4+\mathcal R_{m,\lambda}(\rho),
 \label{new:quartic}\\
 |\mathcal R_{m,\lambda}(\rho)|&\le C|\rho|^6
       \qquad (|\rho|<r).\label{new:remainder}
\end{align}
For negative $\rho$, the deficit in this display denotes the even
extension given by the right side of \eqref{new:exact-deficit}.
All its odd derivatives at zero vanish. In particular there is no cubic
term. The leading coefficient satisfies
\begin{equation}\label{new:curvature-check}
 \frac{\lambda p'(\lambda)}{2m}=\Lambda_\perp
       =\frac{m-1}{2}C_\lambda.
\end{equation}
\end{theorem}

\begin{proof}
All real exponential moments of the Gaussian maximum in \eqref{eq:pc}
are finite, so its moment generating function, and hence $p$, is real
analytic. Taylor's theorem through the quadratic term at $\lambda$
has error $O(|s-\lambda|^3)$.
Substituting $s=\lambda-\lambda\rho^2/(2m)$ gives
\eqref{new:quartic}--\eqref{new:remainder}; the sign of the quartic term
comes from subtracting $p(s)$ from $p(\lambda)$.
The exact expression is even in $\rho$, which gives the assertion about
all odd derivatives. The normal covariance path has initial velocity
$(J-I_m)/m$ in the squared radius, so its first deficit derivative is
$C_\lambda(m-1)/2$. Differentiating the exact value gives the other side
of \eqref{new:curvature-check}. Detailed domination and remainder
arguments are in Appendix~\ref{ap:normal}.
\end{proof}

The equal-curvature boundary is excluded, and no claim of a finite-radius
in-plane optimizer is made. The positive radius and remainder constant
are existential; they do not supply an effective numerical certificate.
The channel crossings, the strict-normal optimization theorem and the
quantified Taylor expansion are formalized over the actual Gaussian
integrals and orthogonal Procrustes distance. Appendix~\ref{ap:lean}
records the exact scope, theorem correspondence and reproducible Lean
checks; the other global bounds and numerical experiments in this paper
are not included in that formalization.


\section{Finite-dimensional decoding diagnostics}\label{sec:application}
The local coefficients describe directional sensitivity near a simplex.
We test that use separately from fixed-budget learning and from the numerical
size of the global recovery constant. All comparisons use the channel in
\eqref{eq:pc}, uniform messages and unit-energy codewords.

\subsection{Directional prediction}
We use $(m,d)=(4,3)$ and $(16,15)$, with $\lambda\in\{1,2,3\}$.
For each of the two in-plane sectors, three independently generated tangent
directions have unit configuration Frobenius norm. Spherical geodesics preserve
each codeword's energy. At amplitudes $t\in\{0.05,0.10,0.20\}$, the target is
the central average $[\delta_\lambda(X(t))+\delta_\lambda(X(-t))]/2$.
The predictor is $\Lambda_jt^2$ for a direction in sector $j$.
Two scalar controls replace $\Lambda_j$ by its dimension-weighted average or
by the equally weighted average of the two sectors.
These controls measure the effect of discarding direction information.

The resulting 108 cells use eight independently randomized Sobol integrations,
with $2^{11}$ points per replicate. Table~\ref{tab:directional} reports mean
absolute prediction error over 18 cells per row.
The scalar controls have the same displayed aggregate error on this balanced
grid, although their cellwise predictions differ.
The directional predictor has smaller aggregate error in every row.
This comparison does not establish a uniform cellwise ordering.
For one four-signal cell at $\lambda=3$ and $t=0.20$, its error is
$2.54\times10^{-5}$, compared with $1.99\times10^{-5}$ for the equal-sector control.

\begin{table}[htbp]
\centering
\caption{Mean absolute error for the central decoding-loss prediction.
Each row aggregates both in-plane sectors, three directions and three channel parameters.}
\label{tab:directional}
\begin{tabular}{rrrrr}
\toprule
$m$ & $t$ & Directional & Dimension-weighted & Equal-sector\\
\midrule
4 & 0.05 & $3.89\times10^{-8}$ & $3.25\times10^{-6}$ & $3.25\times10^{-6}$\\
4 & 0.10 & $6.68\times10^{-7}$ & $1.31\times10^{-5}$ & $1.31\times10^{-5}$\\
4 & 0.20 & $1.09\times10^{-5}$ & $5.43\times10^{-5}$ & $5.43\times10^{-5}$\\
16 & 0.05 & $4.47\times10^{-8}$ & $1.08\times10^{-5}$ & $1.08\times10^{-5}$\\
16 & 0.10 & $1.88\times10^{-7}$ & $4.31\times10^{-5}$ & $4.31\times10^{-5}$\\
16 & 0.20 & $1.01\times10^{-6}$ & $1.73\times10^{-4}$ & $1.73\times10^{-4}$\\
\bottomrule
\end{tabular}
\end{table}

Central averaging removes odd-order terms and changes the prediction target.
For a four-signal direction at $\lambda=2$ and $t=0.20$, the two losses are
$0.0026961$ and $0.0023393$. Their average is $0.0025177$, while the quadratic
prediction is $0.0025202$. A one-sided tolerance requires a separate remainder bound.
All 108 losses exceed five estimated integration standard errors.
That check resolves the loss, rather than the much smaller prediction residual.
Eight prespecified higher-resolution checks have paired difference intervals
containing zero. Their differences can still exceed the residuals in
Table~\ref{tab:directional}. At $m=16$ and $\lambda=2$, one checked direction
changes by $-1.48\times10^{-7}$ at $t=0.05$ and by $-2.37\times10^{-6}$ at $t=0.20$.
The table therefore supports a coarse directional diagnostic, without a
certified residual, improvement factor or higher-order coefficient.

\subsection{A fixed-budget learning comparison}
We adapt the communication autoencoder of O'Shea and Hoydis
\cite[Section III-A and Table IV]{OSheaHoydis} to the two dimensions above.
The encoder uses a dense layer of width $m$ with rectified linear activation,
a linear layer of width $d$ and per-codeword energy normalization.
The decoder uses a dense layer of width $m$ with rectified linear activation
and a softmax output of width $m$.
Training minimizes cross-entropy at $E_b/N_0=7$ dB, where
$E_b/N_0=\lambda^2/(2\log_2m)$.
An additional arm uses the ordinary simplex Gram mean-square penalty with
weight $0.1$. This penalty follows from the known optimal geometry.
The local decoding Hessian is not a cross-entropy Hessian or an Adam preconditioner.

Five paired initializations per dimension use identical message and noise streams.
Both arms use double precision and Adam with learning rate $0.001$ for 1200 steps.
The batch size is 256 for the first 600 steps and 1024 for the remaining 600 steps.
Every run contributes its final checkpoint, with no seed or checkpoint selection.
These 20 runs complete a fixed budget, without an optimization-convergence claim.
The dimensions differ from the original autoencoder's $(16,7)$ example.

Table~\ref{tab:learning} evaluates each learned codebook with matched
maximum-likelihood decoding. Block error rate is one minus decoding success
probability. Each interval is a paired Student-$t$ interval over five training
initializations, evaluated pointwise without multiplicity correction.
The eight mean differences favor the Gram penalty under this budget.
Both arms retain a gap from the ideal simplex.
The learned decoder is a different endpoint. In two of the five 16-signal pairs,
its estimated error increases at all four tested channel parameters.
Cross-initialization paired intervals for that decoder were not evaluated,
so the comparison establishes no end-to-end decoder improvement.

\begin{table}[htbp]
\centering\small
\caption{Mean matched-decoder block error rate over five paired initializations.
The interval is for the Gram-penalized arm minus the baseline.}
\label{tab:learning}
\begin{tabular}{rrrrrl}
\toprule
$m$ & $\lambda$ & Baseline & Gram penalty & Simplex & Paired 95\% interval\\
\midrule
4 & 1 & 0.421808 & 0.403936 & 0.399550 & $[-0.032181,-0.003562]$\\
4 & 2 & 0.151550 & 0.127305 & 0.121023 & $[-0.043128,-0.005362]$\\
4 & 3 & 0.034736 & 0.022434 & 0.019231 & $[-0.021918,-0.002685]$\\
4 & 4.47744 & 0.001919 & 0.000667 & 0.000375 & $[-0.002372,-0.000131]$\\
16 & 1 & 0.749215 & 0.746866 & 0.729857 & $[-0.003146,-0.001552]$\\
16 & 2 & 0.434459 & 0.428081 & 0.383069 & $[-0.008566,-0.004190]$\\
16 & 3 & 0.165418 & 0.158813 & 0.116521 & $[-0.008638,-0.004573]$\\
16 & 6.33206 & 0.0003243 & 0.0002622 & 0.0000273 & $[-0.0000924,-0.0000318]$\\
\bottomrule
\end{tabular}
\end{table}

Evaluation retains independent direct Gaussian sampling with 100000 trials
per codebook and channel parameter. The integration evaluation uses eight
randomized Sobol replicates with $2^{12}$ points each.
Training variability and integration variability remain distinct sources of uncertainty.
The computations used four CPU workers under Linux, Python 3.11.16,
NumPy 2.2.0, SciPy 1.17.1, PyTorch 2.8.0 and mpmath 1.3.0.
No timing comparison is made.

\subsection{Size of the global recovery constant}
The explicit constant in \eqref{eq:codeconstants} is positive but very small.
Table~\ref{tab:global-constant} evaluates a conservative partition envelope
for its scalar maximization. Repeating the $\lambda=2$ cases at 80 and 120
decimal digits gives consistent values. Such agreement is not outward-rounded
interval certification.

\begin{table}[htbp]
\centering
\caption{Numerical values of the explicit global Procrustes constant $c_P$.}
\label{tab:global-constant}
\begin{tabular}{rrrr}
\toprule
$m$ & $\lambda=1$ & $\lambda=2$ & $\lambda=3$\\
\midrule
4 & $2.1293\times10^{-39}$ & $1.2026\times10^{-39}$ & $1.4640\times10^{-40}$\\
16 & $4.6152\times10^{-233}$ & $3.5491\times10^{-233}$ & $7.5743\times10^{-234}$\\
\bottomrule
\end{tabular}
\end{table}

At $\lambda=2$, certifying $d_P\le0.1$ through this inverse bound would
require a reliable deficit upper bound of $1.20\times10^{-41}$ for four signals,
or $3.55\times10^{-235}$ for 16 signals.
These values are far below the precision of the empirical estimates.
The explicit global bound supplies a rate, while practical geometric recovery
requires sharper constants and an effective numerical lower bound for the local transition radius.

\section{Conclusion and limitations}
The explicit comparison controls Gram defect and centroid displacement,
including singular configurations. Together they give a global quadratic
configuration bound at fixed $m$ and positive $\lambda$.
The local expansion identifies the feasible least-loss sectors in the
configuration metric. Numerical diagnostics support its aggregate use for
small central perturbations of four- and 16-signal constellations.
The normalized Gaussian moment family has three unique ordered crossings
for $m\ge4$ and a separate unique crossing for $m=3$.
In the strict-normal regime the finite-radius small-shell optimizer is
exactly the common-normal deformation, up to orthogonal equivalence.
Its deficit has a quartic expansion with a quantified $O(\rho^6)$
remainder and no odd Taylor terms. The radius and remainder constant
depend only on $m$ and $\lambda$, not on the embedding dimension.
The global comparison, recovery and mean-width results, together with
the transition, strict-normal optimization and remainder results, have
been checked in Lean over the actual Gaussian integrals. The global
constants are included; numerical diagnostics and experimental claims
remain outside the certificate.

The global constants deteriorate rapidly with the number of signals.
Their numerical values provide no practical geometry certificate in the
tested cases. Practical recovery needs sharper constants, an effective numerical lower
bound for the local transition radius, computable remainder bounds,
and controlled integration error.
The fixed-budget learning comparison concerns known Gram regularization
and matched decoding. Its decoder and convergence limitations remain separate
from the mathematical stability result.
The rank penalty quantifies a gap from the unrestricted simplex value,
without identifying the optimal rank-constrained constellation.
The channel model assumes equal message probabilities and isotropic Gaussian
noise. Fading, synchronization, channel estimation and constrained modulation
alphabets require separate analysis.

\appendix
\numberwithin{theorem}{section}
\numberwithin{equation}{section}
\section{Proof of the channel transitions}\label{ap:transitions}

This appendix proves the scalar and integral assertions underlying
Theorems~\ref{thm:ordered-switches} and~\ref{thm:triangular-switch}.
Fix an integer \(m\ge3\) and put \(\alpha=(m-1)/m\).
For \(c\in\R\), let
\begin{gather*}
 \phi(c)=(2\pi)^{-1/2}e^{-c^2/2},\qquad
 \Phi(c)=\int_{-\infty}^{c}\phi(s)\,ds,\\
 q(c)=\frac{\phi(c)}{\Phi(c)},\qquad
 h(c)=c+q(c),\qquad \sigma(c)=1-q(c)h(c).
\end{gather*}
These functions are smooth because \(0<\Phi(c)<1\).
For \(\mu\in\R\), define
\begin{equation}\label{sc:tilted-law}
 w_{m,\mu}(c)=e^{\mu c}\phi(c)^2\Phi(c)^{m-2},\qquad
 Z_m(\mu)=\int_{\R}w_{m,\mu}(c)\,dc,\qquad
 p_{m,\mu}(c)=\frac{w_{m,\mu}(c)}{Z_m(\mu)} .
\end{equation}
The density is positive and has Gaussian tails; in particular
\(0<Z_m(\mu)<\infty\). Write
\[
 a_m(\mu)=\E_{m,\mu}[h^2],\qquad
 b_m(\mu)=\E_{m,\mu}[\sigma],\qquad
 H_m(\mu)=\E_{m,\mu}[h].
\]
We omit \(m\) when it is fixed. The channel parameter of
Section~\ref{sec:new-results} is \(\lambda>0\), and
\(\mu=\lambda/\sqrt{\alpha}\).
The endpoint \(\mu=0\) is used only for the continuous family of these
normalized moments. Define
\begin{equation}\label{sc:normalized-curvatures}
 L_0(m,a)=\frac{(m-1)a}{m},\qquad
 L_1(m,a,b)=\frac{a+(m-2)b}{2},\qquad
 L_\perp(m)=\frac{m-1}{2}.
\end{equation}
Their relation to the unnormalized decoding curvatures is recorded in
Subsection~\ref{sc:sub9}.


\subsection{Scalar differential identities}\label{sc:sub1}

\begin{lemma}\label{sc:lemma1.1}
\(q'=-qh\), \(h'=\sigma\), and \(\sigma'=q(h^2-\sigma)\).

\end{lemma}
\begin{proof}
Since \(\phi'=-c\phi\),
\[
q'=\frac{\phi'\Phi-\phi^2}{\Phi^2}=-q\Bigl(c+\frac{\phi}{\Phi}\Bigr)=-qh.
\]
Then \(h'=1+q'=1-qh=\sigma\), and
\[
\sigma'=-q'h-qh'=q h\cdot h-q\sigma=q(h^2-\sigma).
\]

\end{proof}

\begin{lemma}\label{sc:lemma1.2}
For every real \(c\), one has \(h(c)>0\) and \(0<\sigma(c)<1\). In particular \(h\) and \(h^2\) are strictly increasing, and \(q\) is strictly decreasing. Also
\[
0<h(c)\le \max(c,0)+\sqrt{2/\pi},
\]
with equality on the right at \(c=0\).

\end{lemma}
\begin{proof}
Let \(g\) be standard normal and let \(Y\) be the law of \(c+g\) conditioned on \(c+g>0\). The conditioning event has probability \(\Phi(c)\in(0,1)\), and the conditional law is supported on \((0,\infty)\) with a non-degenerate absolutely continuous distribution, so its mean is positive and its variance is positive and finite. Write \(Z=c+g\), so \(Z\) has density \(\phi(z-c)\) and \(P(Z>0)=\Phi(c)\). The substitution \(u=z-c\) and the identity \((\phi)'=-u\phi\) give
\[
\E[Z\mathbf 1_{Z>0}]=\int_{-c}^{\infty}(u+c)\phi(u)\,du=\phi(c)+c\Phi(c),
\]
because \(\int_{-c}^{\infty}u\phi(u)\,du=\phi(c)\). Dividing by \(\Phi(c)\) produces \(\E[Y]=c+q(c)=h(c)\). For the second moment, \(u^2\phi(u)=\phi(u)-(u\phi(u))'\) and \([u\phi(u)]_{-c}^{\infty}=c\phi(c)\), so \(\int_{-c}^{\infty}u^2\phi=\Phi(c)-c\phi(c)\). Adding \(\int_{-c}^{\infty}(2cu+c^2)\phi=2c\phi(c)+c^2\Phi(c)\) yields
\[
\E[Z^2\mathbf 1_{Z>0}]=\Phi(c)(1+c^2)+c\phi(c),
\]
hence \(\E[Y^2]=1+c^2+cq(c)=1+ch(c)\). Therefore \(\mathrm{Var}(Y)=1+ch-h^2\). The identity \(h=c+q\) turns the right-hand side into \(1-qh=\sigma\). Thus \(h>0\) and \(\sigma>0\). Since \(q>0\) and \(h>0\), one has \(qh>0\), so \(\sigma<1\).

Now \(h'=\sigma>0\), so \(h\) is strictly increasing, and \((h^2)'=2h\sigma>0\). Also \(q'=-qh<0\), so \(q\) is strictly decreasing. Therefore \(h(c)<h(0)\) for \(c<0\) and \(q(c)\le q(0)\) for \(c\ge 0\). The value \(h(0)=q(0)=\phi(0)/\Phi(0)=\sqrt{2/\pi}\) gives the stated bound.

\end{proof}

\subsection{The zero-locus derivatives}\label{sc:sub2}

Define the scalar functions
\[
F(c)=h(c)^2-\sigma(c),\qquad D(c)=\sigma+\sigma^2-(1-\sigma)h^2.
\]
Both are smooth. The chain rule and Lemma~\ref{sc:lemma1.1} give, at every \(c\),
\[
F'=2h\sigma-q F.
\]
In particular, on the locus \(F=0\),
\[
F'=2h\sigma>0,
\]
because \(h>0\) and \(\sigma>0\).

For \(D\), use \(q=(1-\sigma)/h\), which is the rearrangement of \(\sigma=1-qh\). Then
\[
\sigma'=\frac{(1-\sigma)(h^2-\sigma)}{h},\qquad (h^2)'=2h\sigma.
\]
Writing \(s=\sigma(c)\) and \(x=h(c)^2\) at a fixed \(c\),
\begin{equation}
\label{sc:Dprime-expr}
D'(c)=\frac{(1-s)(x-s)}{h}\,(1+2s+x)-2hs(1-s).
\end{equation}
This is an identity of smooth functions, obtained by the chain rule, not by differentiating a remainder.

\begin{lemma}[positivity at a zero]\label{sc:lemma2.1}
Let \(h,s\in\R\) satisfy \(0<h\), \(0<s<1\), and
\[
s+s^2-(1-s)h^2=0.
\]
Then \(x=h^2\) equals \(s(1+s)/(1-s)\), and the right-hand side of \eqref{sc:Dprime-expr} equals
\[
\frac{4s^3}{h(1-s)}>0.
\]
Thus, whenever \(D(c)=0\), one has \(D'(c)>0\).

\end{lemma}
\begin{proof}
The constraint solves as \(x=s(1+s)/(1-s)\). Then
\[
x-s=\frac{2s^2}{1-s},\qquad \sigma'\text{ coefficient}=\frac{1-s}{h}\cdot\frac{2s^2}{1-s}=\frac{2s^2}{h}.
\]
Substitute \(x\) into \(1+2s+x\):
\[
1+2s+x=\frac{(1+2s)(1-s)+s(1+s)}{1-s}=\frac{1+2s-s^2}{1-s}.
\]
The expression \eqref{sc:Dprime-expr} becomes
\begin{align*}
&\frac{2s^2}{h}\cdot\frac{1+2s-s^2}{1-s}-2hs(1-s)\\
&\qquad=\frac{2s^2(1+2s-s^2)-2h^2 s(1-s)^2}{h(1-s)}.
\end{align*}
The relation \(h^2=s(1+s)/(1-s)\) gives \(h^2(1-s)^2=s(1+s)(1-s)=s(1-s^2)\), so the second numerator term is \(2s^2(1-s^2)\). The numerator collapses to
\[
2s^2\bigl(1+2s-s^2-1+s^2\bigr)=4s^3.
\]
Hence the expression equals \(4s^3/(h(1-s))\). The hypotheses \(h>0\) and \(0<s<1\) make this positive.

\end{proof}


The same computation with no constraint shows the sign identity used later for the ratio:
\[
\Bigl(\frac{h^2}{\sigma}\Bigr)'=\frac{h\,D}{\sigma^2}.
\]
Indeed, with \(x=h^2\) and \(s=\sigma\),
\[
\Bigl(\frac{x}{s}\Bigr)'=\frac{x's-xs'}{s^2},\qquad x'=2hs,\qquad s'=\frac{(1-s)(x-s)}{h},
\]
and clearing the positive factor \(h\) produces the numerator \(hD\) if and only if
\[
2s^2-(1-s)(x-s)=s+s^2-(1-s)x,
\]
which is an algebraic identity. Thus, wherever \(D>0\), the ratio \(h^2/\sigma\) is strictly increasing. The sign of \(D\) is proved in Subsection~\ref{sc:sub4}, after the tail expansion.

\subsection{Mills expansion with a bounded remainder}\label{sc:sub3}

For \(t>0\) the left-tail Mills ratio is
\[
M(t)=\frac{\Phi(-t)}{\phi(t)}=\int_0^\infty\exp(-tv-v^2/2)\,dv.
\]
The substitution \(s=tv\) gives the Laplace representation,
\[
M(t)=t^{-1}\int_0^\infty e^{-s}\exp\bigl(-s^2/(2t^2)\bigr)\,ds.
\]
At \(c=-t\) one has \(q(-t)=1/M(t)\) and \(h(-t)=1/M(t)-t\).

\begin{lemma}[explicit Taylor remainder]\label{sc:lemma3.1}
Let \(n\ge1\) be an integer and let \(t>0\). Then
\[
M(t)=\sum_{k=0}^{n-1}(-1)^k\frac{(2k)!}{k!\,2^k}\,t^{-(2k+1)}+\theta_n(t),
\]
where
\[
|\theta_n(t)|\le\frac{(2n)!}{n!\,2^n\,t^{2n+1}}.
\]
In particular, for \(n=5\),
\[
M(t)=t^{-1}-t^{-3}+3t^{-5}-15t^{-7}+105t^{-9}+\theta_5(t),
\]
with \(|\theta_5(t)|\le 945\,t^{-11}\). The remainder is not differentiated anywhere below.

\end{lemma}
\begin{proof}
For \(x>0\) the Lagrange remainder after the terms of degree \(0,\ldots,n-1\) in \(e^{-x}\) has absolute value at most \(x^n/n!\), because it equals \((-1)^n e^{-\xi}x^n/n!\) for some \(\xi\in(0,x)\) and \(e^{-\xi}\le1\); at \(x=0\) the remainder is zero. Insert \(x=s^2/(2t^2)\), integrate against \(e^{-s}\) on \((0,\infty)\), and use \(\int_0^\infty e^{-s}s^{2n}\,ds=(2n)!\). Dividing by \(t\) gives the bound on \(\theta_n\). The coefficient of \(t^{-(2k+1)}\) is \((-1)^k(2k)!/(k!2^k)\). For \(k\le 4\) these coefficients are \(1,-1,3,-15,105\), and the \(n=5\) bound is \(10!/(120\cdot 32)=945\).

\end{proof}


Set \(u=t^{-2}\). For every \(t>0\),
\begin{equation}
\label{sc:M-factor}
M(t)=t^{-1}\bigl(1-u+3u^2-15u^3+105u^4+r(t)\bigr),
\end{equation}
where \(|r(t)|\le 945\,u^5\). Write \(a_0(u)=-u+3u^2-15u^3+105u^4\) and \(a(t)=a_0(u)+r(t)\), so the factor in parentheses is \(1+a(t)\).

\begin{lemma}[algebraic inversion on a far ray]\label{sc:lemma3.2}
There exist a finite \(T_0\ge 1\) and a finite positive constant \(C_h\) such that for every \(t\ge T_0\),
\begin{equation}
\label{sc:h-exp}
h(-t)=t^{-1}-2t^{-3}+10t^{-5}-74t^{-7}+\rho_h(t),
\qquad |\rho_h(t)|\le C_h\,t^{-9}.
\end{equation}
The coefficient \(-74\) is required later. Stopping the expansion of \(h\) at \(O(t^{-7})\) does not determine the coefficient of \(t^{-6}\) in \(\sigma\).

\end{lemma}
\begin{proof}
As \(u\to 0\), \(a_0(u)\to 0\) and \(r(t)\to 0\). Choose a finite \(T_0\ge 1\) such that \(t\ge T_0\) implies both \(|a_0(u)|\le 1/4\) and \(|r(t)|\le 1/4\). Then \(|a(t)|\le 1/2\) and \(|a_0(u)|\le 1/2\), so \(1+a(t)\ge 1/2\) and \(q(-t)=t/(1+a(t))\) is well-defined. In particular
\[
h(-t)=t\Bigl(\frac{1}{1+a(t)}-1\Bigr).
\]
On the compact interval \(0\le u\le T_0^{-2}\) one has \(|a_0(u)|\le 1/2\), so \(g(u)=(1+a_0(u))^{-1}\) is smooth there. The recurrence below computes the Taylor coefficients of \(g\) at \(0\), because it is the Cauchy product for the reciprocal of the polynomial \(1+a_0\). Lagrange's theorem in the variable \(u\), through degree \(4\), then gives a remainder bounded by a constant times \(u^5\): the fifth derivative of \(g\) is continuous on that compact interval. The recurrence \(b_0=1\) and \(b_n=-\sum_{j=1}^n a_j b_{n-j}\) for the coefficients of \((1+a_0)^{-1}\), with \(a_1=-1\), \(a_2=3\), \(a_3=-15\), \(a_4=105\), gives
\begin{align*}
b_1&=1,\\
b_2&=-(-1+3)=-2,\\
b_3&=-(2+3-15)=10,\\
b_4&=-(-10-6-15+105)=-74.
\end{align*}
Thus
\[
(1+a_0)^{-1}=1+u-2u^2+10u^3-74u^4+O(u^5),
\]
where the \(O(u^5)\) on \(t\ge T_0\) is the Lagrange remainder of \(g\). Both \(a_0(u)\) and \(a(t)\) lie in \([-1/2,1/2]\), and so does the segment between them. On that segment the derivative of \((1+z)^{-1}\) has absolute value at most \(4\), so the mean value theorem gives
\[
\bigl|(1+a)^{-1}-(1+a_0)^{-1}\bigr|\le 4|r(t)|\le 3780\,u^5.
\]
Thus the same expansion holds for \((1+a)^{-1}\), with a possibly larger finite constant. Multiplying by \(t\) and subtracting \(t\) produces \eqref{sc:h-exp}. No derivative of \(\rho_h\) is taken.

\end{proof}

\begin{lemma}\label{sc:lemma3.3}
There exist a finite \(T_1\ge T_0\) and finite positive constants \(C_\sigma,C_F,C_D\) such that for every \(t\ge T_1\),
\begin{align*}
\sigma(-t)&=t^{-2}-6t^{-4}+50t^{-6}+\rho_\sigma(t),
&|\rho_\sigma|&\le C_\sigma t^{-8},\\[4pt]
F(-t)&=2t^{-4}+\rho_F(t),
&|\rho_F|&\le C_F t^{-6},\\[4pt]
D(-t)&=4t^{-6}+\rho_D(t),
&|\rho_D|&\le C_D t^{-8}.
\end{align*}
In particular there exists a finite \(T_*\) such that \(F(-t)>0\) and \(D(-t)>0\) for all \(t\ge T_*\).

\end{lemma}
\begin{proof}
Retain \(t\ge T_0\) and write \(H=th(-t)\), so Lemma~\ref{sc:lemma3.2} says
\[
H=1-2u+10u^2-74u^3+\tilde\rho,\qquad |\tilde\rho|\le C_h u^4.
\]
Set \(H_1=H-1=-2u+10u^2-74u^3+\tilde\rho\). From \(q=t+h=t(1+uH)\) and \(\sigma=1-qh\),
\[
qh=H+uH^2,\qquad \sigma=1-H-uH^2=-H_1-u-2uH_1-uH_1^2.
\]
Expand, using only products of series and the bound on \(\tilde\rho\). The terms through \(u^3\) are
\begin{align*}
-H_1&=2u-10u^2+74u^3+O(u^4),\\
-2uH_1&=4u^2-20u^3+O(u^4),\\
H_1^2&=4u^2-40u^3+O(u^4),\\
-uH_1^2&=-4u^3+O(u^4).
\end{align*}
Adding and using \(-u\) gives
\[
\sigma=u-6u^2+50u^3+O(u^4).
\]
The summand \(74u^3\) comes from the coefficient \(-74\) in \(h\). It is the reason that coefficient must be kept. The symbol \(O(u^4)\) denotes a remainder bounded by a constant times \(u^4\) for \(t\ge T_0\); the constant exists because it is a finite combination of \(C_h\) and the absolute values of the displayed coefficients. This is the expansion of \(\sigma\).

Next, \(h^2=uH^2=u(1+2H_1+H_1^2)\). Through order \(u^3\),
\[
h^2=u-4u^2+24u^3+O(u^4).
\]
The coefficient \(24\) uses \(h\) only through \(t^{-5}\). Subtracting the expansion of \(\sigma\),
\[
F=h^2-\sigma=2u^2+O(u^3)=2t^{-4}+O(t^{-6}).
\]
The leading coefficient \(2\) requires \(h\) through its \(10t^{-5}\) term, which fixes the coefficient \(-6t^{-4}\) in \(\sigma\). Choose the constant \(C_F\) from this \(O(t^{-6})\).

For \(D\), square the expansion of \(\sigma\):
\[
\sigma^2=u^2-12u^3+O(u^4),
\]
so \(\sigma+\sigma^2=u-5u^2+38u^3+O(u^4)\). Also
\[
(1-\sigma)h^2=u-5u^2+34u^3+O(u^4),
\]
where the coefficient \(34\) is \(24+4+6\), coming from \(h^2\), from \((-\,u)(-4u^2)\), and from \((6u^2)(u)\). Therefore
\[
D=4u^3+O(u^4)=4t^{-6}+O(t^{-8}).
\]
The leading coefficient \(4\) uses the coefficient \(50\) in \(\sigma\), hence uses the coefficient \(-74\) in \(h\). The constant \(C_D\) exists for the same reason as \(C_F\).

Finally \(F(-t)\ge t^{-4}(2-C_F t^{-2})\) for \(t\ge T_1\), and \(D(-t)\ge t^{-6}(4-C_D t^{-2})\) for \(t\ge T_1\). Any
\[
T_*>\max\bigl(T_1,\sqrt{C_F},\sqrt{C_D/4}\bigr)
\]
is a finite cutoff on which both functions are positive. The value of \(T_*\) is not required.

\end{proof}

\subsection{First-zero argument}\label{sc:sub4}

\begin{lemma}\label{sc:lemma4.1}
Let \(f\in C^1(\R)\). Suppose there is a finite \(c_0\) such that \(f>0\) on \((-\infty,c_0]\), and suppose \(f'(c)>0\) whenever \(f(c)=0\). Then \(f>0\) on \(\R\).

\end{lemma}
\begin{proof}
If the set \(E=\{c:f(c)\le 0\}\) is empty, there is nothing to prove. Otherwise \(c_*=\inf E\) is finite and \(c_*>c_0\). Continuity gives \(f(c_*)=0\), while \(f>0\) on \((-\infty,c_*)\). Differentiability gives
\[
f(c_*-\varepsilon)=-\varepsilon f'(c_*)+o(\varepsilon).
\]
The hypothesis \(f'(c_*)>0\) makes the right-hand side negative for small \(\varepsilon>0\), contradicting positivity on the left. The remainder in this one-variable Taylor expansion is the definition of \(f'(c_*)\); it is not an asymptotic remainder of \(M(t)\).

\end{proof}

\begin{theorem}\label{sc:theorem4.2}
The functions \(F\) and \(D\) are positive on all of \(\R\). Consequently \(\sigma'=qF>0\) everywhere, so \(\sigma\) is strictly increasing, and \((h^2/\sigma)'=hD/\sigma^2>0\) everywhere, so \(h^2/\sigma\) is strictly increasing.

\end{theorem}
\begin{proof}
Lemma~\ref{sc:lemma3.3} supplies a finite left endpoint beyond which \(F(-t)>0\) and \(D(-t)>0\). On every zero of \(F\), Subsection~\ref{sc:sub2} gives \(F'=2h\sigma>0\). On every zero of \(D\), Lemma~\ref{sc:lemma2.1} gives \(D'>0\). Lemma~\ref{sc:lemma4.1} applies to both functions. The sign statements for the derivatives of \(\sigma\) and \(h^2/\sigma\) were proved in Subsections~\ref{sc:sub1} and~\ref{sc:sub2}, conditionally on \(F>0\) and \(D>0\).

Combined with Lemma~\ref{sc:lemma1.2}, the four scalar maps \(h\), \(h^2\), \(\sigma\), and \(h^2/\sigma\) are strictly increasing on \(\R\).

\end{proof}

\subsection{Differentiation under the integral and strict covariance}\label{sc:sub5}

Let \(w_{m,\mu}(c)=e^{\mu c}\phi(c)^2\Phi(c)^{m-2}\) and let \(Z_m(\mu)=\int_{\R}w_{m,\mu}\). For each fixed \(m\) the map \(\mu\mapsto w_{m,\mu}(c)\) is smooth. For \(|\mu|\le M\), the estimate
\[
 |c|^N w_{m,\mu}(c)
 \le \frac{e^{M^2/2}}{2\pi}(1+|c|)^N e^{-c^2/2}
\]
follows from \(M|c|\le(c^2+M^2)/2\). It gives a common integrable bound for every fixed polynomial degree.

\begin{lemma}\label{sc:lemma5.1}
For every integer \(m\ge 3\) and every real \(\mu\), the normalizing constant \(Z_m\) is smooth and positive, and if \(f:\R\to\R\) is continuous and \(|f(c)|\le C(1+|c|)^N\), then \(\mu\mapsto\E_{m,\mu}[f]\) is smooth and
\[
\frac{d}{d\mu}\E_{m,\mu}[f(c)]=\mathrm{Cov}_{m,\mu}\bigl(f(c),c\bigr).
\]
If in addition \(f\) is strictly increasing, the covariance is strictly positive.

\end{lemma}
\begin{proof}
The difference quotient of \(w_{m,\mu}\) in \(\mu\) is \(w_{m,\mu}(c)\cdot c\cdot e^{\theta c}/e^{\mu c}\) for an intermediate \(\theta\), hence is dominated on a neighborhood of any fixed \(\mu\) by an integrable function of the type above. Differentiation under the integral sign is therefore legitimate and gives \(Z_m'(\mu)=\int c\,w_{m,\mu}\). The same domination applies to \(f(c)w_{m,\mu}(c)\) and to \(c f(c)w_{m,\mu}(c)\), so the quotient defining the expectation may be differentiated by the quotient rule. The result is the covariance.

For the strict sign, let \(X,Y\) be independent copies of the law \(p_{m,\mu}\). Then
\[
\mathrm{Cov}(f(X),X)=\frac12\E\bigl[(f(X)-f(Y))(X-Y)\bigr].
\]
The integrand is nonnegative. It is not almost surely zero: the density \(p_{m,\mu}\) is continuous and strictly positive on \(\R\), so \(P(X\neq Y)=1\), and a strictly increasing \(f\) then gives a strictly positive integrand on \(\{X\neq Y\}\). The expectations exist by the polynomial-growth hypothesis and the Gaussian tail of \(p_{m,\mu}\).

\end{proof}

\begin{lemma}\label{sc:lemma5.2}
For each integer \(m\ge 3\), the maps \(a_m\), \(b_m\), \(H_m\), and \(a_m/b_m\) are smooth and strictly increasing on \(\R\). In particular they are strictly increasing on \([0,\infty)\), and \(H_m(\mu)\ge H_m(0)>0\) for \(\mu\ge 0\).

\end{lemma}
\begin{proof}
The functions \(h^2\), \(\sigma\), and \(h\) are continuous, strictly increasing, and of polynomial growth by Lemma~\ref{sc:lemma1.2} and Theorem~\ref{sc:theorem4.2}. Lemma~\ref{sc:lemma5.1} gives the result for \(a_m\), \(b_m\), and \(H_m\). Strict increase of \(H_m\) yields \(H_m(\mu)\ge H_m(0)\). Positivity of \(h\) yields \(H_m(0)>0\).

For the ratio, let \(Q_{m,\mu}\) be the probability with density proportional to \(\sigma(c)w_{m,\mu}(c)\). The factor \(\sigma(c)\) does not depend on \(\mu\), lies in \((0,1)\), and is continuous and positive, so \(Q_{m,\mu}\) is again a smooth exponential family in \(\mu\) with the same dominating Gaussian tails. Moreover
\[
\E_{Q_{m,\mu}}\Bigl[\frac{h^2}{\sigma}\Bigr]=\frac{\E_{m,\mu}[h^2]}{\E_{m,\mu}[\sigma]}=\frac{a_m(\mu)}{b_m(\mu)}.
\]
Theorem~\ref{sc:theorem4.2} says that \(h^2/\sigma\) is strictly increasing. Lemma~\ref{sc:lemma3.3} gives \(h(-t)\sim t^{-1}\) and \(\sigma(-t)\sim t^{-2}\), so \(h^2/\sigma\to 1\) as \(c\to-\infty\). For \(c\ge 0\), the same theorem gives \(\sigma(c)\ge\sigma(0)=1-2/\pi>0\), and Lemma~\ref{sc:lemma1.2} gives \(h(c)\le c+\sqrt{2/\pi}\). Hence
\[
0<\frac{h(c)^2}{\sigma(c)}\le\frac{(c+\sqrt{2/\pi})^2}{1-2/\pi}
\]
on \([0,\infty)\). Continuity controls the remaining compact interval, so the ratio has polynomial growth on all of \(\R\). The weight of \(Q_{m,\mu}\) is \(\sigma(c)w_{m,\mu}(c)\). Since \(0<\sigma<1\), every integrable dominant used for \(w_{m,\mu}\) dominates the corresponding integrand for \(Q_{m,\mu}\), and differentiation in \(\mu\) still multiplies that weight by \(c\). The covariance argument of Lemma~\ref{sc:lemma5.1} therefore applies to \(Q_{m,\mu}\) and gives strict increase in \(\mu\).

\end{proof}

\begin{lemma}[monotonicity in the number of signals]\label{sc:lemma5.3}
Fix \(\mu\in\R\). For every integer \(m\ge 3\),
\[
b_{m+1}(\mu)>b_m(\mu).
\]
The same strict increase holds for \(a_m(\mu)\), \(H_m(\mu)\), and \(a_m(\mu)/b_m(\mu)\).

\end{lemma}
\begin{proof}
The un-normalized weight at \(m+1\) is the weight at \(m\) multiplied by \(\Phi(c)\). The function \(\Phi\) is strictly increasing, and \(\sigma\) is strictly increasing. For a positive probability \(P\) with full support on \(\R\) and a bounded, positive, strictly increasing \(g\), the reweighted expectation of a strictly increasing integrable \(f\) satisfies
\[
\frac{\E[fg]}{\E[g]}-\E[f]=\frac{1}{2\E[g]}\E\bigl[(f(X)-f(Y))(g(X)-g(Y))\bigr]>0,
\]
by the same independent-copy argument as in Lemma~\ref{sc:lemma5.1}. Apply this to \(P=p_{m,\mu}\), \(g=\Phi\), and \(f=\sigma\). Taking \(f=h^2\) or \(f=h\) proves the corresponding assertions for \(a_m\) and \(H_m\). The ratio \(a/b\) is the expectation of the strictly increasing function \(h^2/\sigma\) under the \(\sigma\)-reweighted law, and increasing \(m\) multiplies that law by another strictly increasing copy of \(\Phi\). The polynomial-growth integrability is the same Gaussian estimate as before.

\end{proof}

\subsection{The four-signal benchmark}\label{sc:sub6}

Write \(\mathcal I_4=\int_{\R}\phi^4\) and \(\mathcal J=\int_{\R}\phi^2\Phi^2\).

\begin{lemma}\label{sc:lemma6.1}
\(\mathcal I_4=1/(4\sqrt 2\,\pi^{3/2})\) and
\[
b_4(0)=1-\frac{4\mathcal I_4}{3\mathcal J}.
\]

\end{lemma}
\begin{proof}
The identity \(\phi(c)^4=(2\pi)^{-2}e^{-2c^2}\) and \(\int e^{-2c^2}\,dc=\sqrt{\pi/2}\) give
\[
\mathcal I_4=(2\pi)^{-2}\sqrt{\pi/2}=1/(4\sqrt 2\,\pi^{3/2}).
\]
For the expectation, \(\sigma=1-qh\) and \(qh\,\phi^2\Phi^2=c\phi^3\Phi+\phi^4\). Integrate \(c\phi^3\Phi\) by parts. From \((\phi^3)'=-3c\phi^3\) and the vanishing of \(\phi^3\Phi\) at both ends,
\[
\int c\phi^3\Phi=\frac13\int\phi^4.
\]
Hence \(\int qh\,\phi^2\Phi^2=(4/3)\mathcal I_4\), and dividing by \(\mathcal J\) produces the formula. This identity is exact; the arcsin evaluation of \(\mathcal J\) is not used.

\end{proof}

\begin{lemma}\label{sc:lemma6.2}
\(\mathcal J\ge (3\pi+2)/(24\pi^{3/2})\), and therefore
\[
b_4(0)\ge 1-\frac{4\sqrt 2}{3\pi+2}.
\]

\end{lemma}
\begin{proof}
Expand \(\Phi^2=(\Phi-1/2)^2+(\Phi-1/2)+1/4\). The cross term \(\phi^2(\Phi-1/2)\) is an odd integrable function, so its integral vanishes. Also \(\int\phi^2=1/(2\sqrt\pi)\), hence
\[
\mathcal J=\frac14\int\phi^2+\int\phi^2(\Phi-1/2)^2=\frac{1}{8\sqrt\pi}+\int\phi^2(\Phi-1/2)^2.
\]
Integration by parts, using \((\phi^2)'=-2c\phi^2\) and vanishing boundary terms, gives
\[
\int c\phi^2(\Phi-1/2)=\frac12\int\phi^3.
\]
Cauchy--Schwarz in the form
\[
\Bigl(\int (c\phi)\cdot\bigl(\phi(\Phi-1/2)\bigr)\Bigr)^2
\le\Bigl(\int c^2\phi^2\Bigr)\Bigl(\int\phi^2(\Phi-1/2)^2\Bigr)
\]
therefore yields a lower bound on the second integral. The Gaussian integrals
\[
\int\phi^3=\frac{1}{2\pi\sqrt 3},\qquad \int c^2\phi^2=\frac{1}{4\sqrt\pi}
\]
are elementary. The resulting quotient is \(1/(12\pi^{3/2})\), and
\[
\mathcal J\ge\frac{1}{8\sqrt\pi}+\frac{1}{12\pi^{3/2}}=\frac{3\pi+2}{24\pi^{3/2}}.
\]
Cauchy--Schwarz is not asserted to be equality. Because \(\mathcal I_4>0\), replacing \(\mathcal J\) by this lower bound increases \(4\mathcal I_4/(3\mathcal J)\). The arithmetic
\[
\frac{4\mathcal I_4}{3}\Big/\frac{3\pi+2}{24\pi^{3/2}}=\frac{4\sqrt 2}{3\pi+2}
\]
is the comparison of Lemma~\ref{sc:lemma6.1} with Lemma~\ref{sc:lemma6.2}'s bound.

\end{proof}

\begin{lemma}\label{sc:lemma6.3}
\(\pi>28/9\) and \(\sqrt 2<17/12\). Consequently \(4\sqrt 2/(3\pi+2)<1/2\), and \(b_4(0)>1/2\).

\end{lemma}
\begin{proof}
Let \(\theta=\arctan(1/5)\). Then \(0<\theta<\pi/4\), because \(0<1/5<1\) and tangent is increasing on \((0,\pi/2)\). The double-angle formula applies and gives
\[
\tan(2\theta)=\frac{2\cdot(1/5)}{1-(1/5)^2}=\frac{5}{12}.
\]
Thus \(2\theta\in(0,\pi/2)\). Since \(\tan(2\theta)<1\), one has \(2\theta\in(0,\pi/4)\), and the formula applies again:
\[
\tan(4\theta)=\frac{2\cdot(5/12)}{1-(5/12)^2}=\frac{120}{119},
\]
with \(4\theta\in(0,\pi/2)\). Set \(\eta=4\theta\) and \(\beta=\arctan(1/239)\in(0,\pi/4)\). Then \(\tan\eta=120/119>1/239=\tan\beta\). Both angles lie in \((0,\pi/2)\) and tangent is increasing there, so \(\eta>\beta\) and \(\eta-\beta\in(0,\pi/2)\). Therefore
\begin{align*}
\tan(\eta-\beta)
&=\frac{120/119-1/239}{1+(120/119)(1/239)}
=\frac{120\cdot 239-119}{119\cdot 239+120}\\
&=\frac{28680-119}{28441+120}
=\frac{28561}{28561}=1.
\end{align*}
The unique angle in \((0,\pi/2)\) with tangent \(1\) is \(\pi/4\), so
\[
\frac\pi 4=4\arctan\frac15-\arctan\frac1{239}.
\]
For every real \(t\), the identity \((1-t^2)(1+t^2)+t^4=1\) rearranges to
\[
\frac{1}{1+t^2}=1-t^2+\frac{t^4}{1+t^2}.
\]
Integrating from \(0\) to \(x>0\) gives
\[
\arctan x=x-\frac{x^3}{3}+\int_0^x\frac{t^4}{1+t^2}\,dt>x-\frac{x^3}{3},
\]
the integrand being positive. In particular
\[
\arctan\frac15>\frac15-\frac{1}{3\cdot 5^3}=\frac{74}{375}.
\]
The same representation gives \(\arctan y<y\) for \(y>0\), because \(1/(1+t^2)<1\). Thus \(\arctan(1/239)<1/239\). Combining these strict bounds with Machin's formula,
\[
\frac\pi 4>4\cdot\frac{74}{375}-\frac{1}{239}=\frac{296}{375}-\frac{1}{239}=\frac{70369}{89625}.
\]
The cross multiplication \(70369\cdot 9-7\cdot 89625=5946>0\) gives \(70369/89625>7/9\), so \(\pi>28/9\). Both arctangent inequalities are strict.

Separately, \((17/12)^2=289/144>288/144=2\), so \(\sqrt 2<17/12\). Hence
\[
8\sqrt 2<8\cdot\frac{17}{12}=\frac{34}{3},\qquad 3\pi+2>3\cdot\frac{28}{9}+2=\frac{34}{3}.
\]
The second inequality is strict, so \(8\sqrt 2<3\pi+2\) and \(4\sqrt 2/(3\pi+2)<1/2\). Lemma~\ref{sc:lemma6.2} gives \(b_4(0)\ge 1-4\sqrt 2/(3\pi+2)>1/2\).

\end{proof}

\begin{corollary}\label{sc:corollary6.4}
For every integer \(m\ge 4\) and every \(\mu\ge 0\),
\[
\frac12<b_m(\mu)<1.
\]

\end{corollary}
\begin{proof}
The upper bound holds for every \(m\) and \(\mu\) because \(0<\sigma<1\) pointwise and \(p_{m,\mu}\) has full support, so the expectation of \(\sigma\) cannot attain \(1\). The lower bound holds at \((m,\mu)=(4,0)\) by Lemma~\ref{sc:lemma6.3}. Lemmas~\ref{sc:lemma5.2} and~\ref{sc:lemma5.3} move the value strictly upward in each variable separately, and the inequality is already strict at the minimal corner.

\end{proof}

\subsection{The balance identity and the initial values}\label{sc:sub7}

\begin{lemma}\label{sc:lemma7.1}
For every integer \(m\ge 3\) and every real \(\mu\),
\begin{equation}
\label{sc:identity}
2a_m(\mu)+(m-1)b_m(\mu)=m+\mu H_m(\mu).
\end{equation}

\end{lemma}
\begin{proof}
The logarithmic derivative of the un-normalized weight is
\[
\mu-2c+(m-2)q=\mu-2h+mq,
\]
which is the weight's formula \(p'/p=\mu-2h+mq\). Fix \(R>0\). Then
\[
\int_{-R}^{R}(hp_{m,\mu})'=\bigl[hp_{m,\mu}\bigr]_{-R}^{R}.
\]
The bound of Lemma~\ref{sc:lemma1.2} and the Gaussian tail of \(p_{m,\mu}\) force the boundary term to \(0\) as \(R\to\infty\). Expanding the derivative,
\[
(hp)'=\sigma p+h(\mu-2h+mq)p.
\]
The resulting integrand is absolutely dominated by an integrable multiple of \((1+c^2)p_{m,\mu}(c)\), because \(0<\sigma<1\), \(qh=1-\sigma\), and \(h\) has linear growth. Passing to the limit,
\[
b+\mu H-2a+m\E[qh]=0.
\]
The identity \(\E[qh]=1-b\) rearranges into \eqref{sc:identity}.

\end{proof}

\begin{lemma}\label{sc:lemma7.2}
Fix an integer \(m\ge 4\). At \(\mu=0\),
\[
a(0)<\frac m2,\qquad \frac{a(0)}{b(0)}<m,\qquad a(0)+(m-2)b(0)<m-1.
\]
All three inequalities are strict.

\end{lemma}
\begin{proof}
The identity at \(\mu=0\) is \(2a+(m-1)b=m\), so \(a=[m-(m-1)b]/2\). Since \(b>0\), one has \(a<m/2\). Next
\[
\frac a b=\frac{m}{2b}-\frac{m-1}{2}.
\]
The inequality \(a/b<m\) rearranges to \(b>m/(3m-1)\). For \(m>1\), \(m/(3m-1)<1/2\), and Corollary~\ref{sc:corollary6.4} supplies \(b>1/2\). Finally
\[
a+(m-2)b=\frac{m+(m-3)b}{2}.
\]
Comparing with \(m-1\) rearranges to \((m-3)b<m-2\). For \(m=4\) the right side is \(2\) and the left side is \(b<1\). For \(m>4\), \((m-2)/(m-3)=1+1/(m-3)>1>b\). Every comparison is strict.

\end{proof}

\subsection{Growth and the ordered crossings}\label{sc:sub8}

\begin{lemma}\label{sc:lemma8.1}
For each integer \(m\ge 3\) and every \(\mu\ge 0\),
\[
2a_m(\mu)>1+\mu H_m(0).
\]
In particular \(a_m(\mu)\to+\infty\) as \(\mu\to+\infty\). For \(m=3\), \(\mu H_3(\mu)\ge \mu H_3(0)\), so \(\mu H_3(\mu)\to+\infty\).

\end{lemma}
\begin{proof}
By \eqref{sc:identity} and \(b_m(\mu)<1\),
\begin{align*}
2a_m(\mu)&=m+\mu H_m(\mu)-(m-1)b_m(\mu)\\
&>m+\mu H_m(\mu)-(m-1)=1+\mu H_m(\mu).
\end{align*}
Lemma~\ref{sc:lemma5.2} gives \(H_m(\mu)\ge H_m(0)>0\), and the displayed bound follows. The right-hand side tends to infinity. For \(m=3\) the same lemma gives the linear lower bound \(\mu H_3(0)\).

\end{proof}

\begin{lemma}[comparison of the three coefficients]\label{sc:lemma8.2}
Let \(m\ge 4\), \(a>0\), and \(1/2<b<1\). Then
\begin{align*}
a=\frac m2&\Longrightarrow L_0(m,a)=L_\perp(m)\ \text{and}\ L_\perp(m)<L_1(m,a,b),\\[4pt]
a=mb&\Longrightarrow L_0(m,a)=L_1(m,a,b)\ \text{and}\ L_\perp(m)<L_0(m,a),\\[4pt]
\min\bigl(L_0(m,a),L_\perp(m)\bigr)&<L_1(m,a,b).
\end{align*}

\end{lemma}
\begin{proof}
If \(a=m/2\), then \(L_0=(m-1)/2=L_\perp\) and
\[
L_1-L_\perp=\frac{m-2}{2}\Bigl(b-\frac12\Bigr)>0.
\]
If \(a=mb\), then \(L_1=(mb+(m-2)b)/2=(m-1)b=L_0\), while
\[
L_\perp-L_0=(m-1)\Bigl(\frac12-b\Bigr)<0.
\]
For the third claim, the same algebra gives the unrestricted identity
\[
L_1-L_0=\frac{m-2}{2m}(mb-a)\qquad(m>2).
\]
Suppose toward a contradiction that \(L_1\le L_0\) and \(L_1\le L_\perp\). The first forces \(a\ge mb\). The second forces \(a+(m-2)b\le m-1\). Adding these,
\[
a+(m-2)b\ge mb+(m-2)b=2(m-1)b>(m-1),
\]
which contradicts the second inequality. The last step is strict because \(b>1/2\).

\end{proof}

\begin{theorem}[unique positive crossings for \(m\ge 4\)]\label{sc:theorem8.3}
Fix an integer \(m\ge 4\). There exist unique real numbers
\[
0<\mu_{1\perp}(m)<\mu_{0\perp}(m)<\mu_{01}(m)<\infty
\]
such that
\begin{align*}
a(\mu_{1\perp})+(m-2)b(\mu_{1\perp})&=m-1,\\
a(\mu_{0\perp})&=m/2,\\
a(\mu_{01})&=m\,b(\mu_{01}),
\end{align*}
and each equation has no other root in \([0,\infty)\). Equivalently, the three normalized curvatures satisfy
\[
L_1=L_\perp,\qquad L_0=L_\perp,\qquad L_0=L_1
\]
at these three parameters, in that order.

\end{theorem}
\begin{proof}
The maps
\[
f_{1\perp}=a+(m-2)b,\qquad f_{0\perp}=a,\qquad f_{01}=a/b
\]
are continuous and strictly increasing on \([0,\infty)\) by Lemma~\ref{sc:lemma5.2}. Lemma~\ref{sc:lemma7.2} says that their values at \(0\) are strictly less than \(m-1\), \(m/2\), and \(m\) respectively. Lemma~\ref{sc:lemma8.1} and the bound \(1/2<b<1\) give
\[
f_{0\perp}(\mu)\to\infty,\qquad f_{01}(\mu)\ge a(\mu)\to\infty,\qquad f_{1\perp}(\mu)\ge a(\mu)\to\infty.
\]
The intermediate-value theorem supplies at least one positive root of each equation, and strict increase supplies uniqueness.

Let \(\mu_{0\perp}\) be the unique root of \(a=m/2\). Corollary~\ref{sc:corollary6.4} gives \(b(\mu_{0\perp})>1/2\), so \(a(\mu_{0\perp})<m b(\mu_{0\perp})\). Thus \(f_{01}(\mu_{0\perp})<m\), and the unique root \(\mu_{01}\) of \(f_{01}=m\) lies strictly to the right of \(\mu_{0\perp}\). At the same point,
\[
a+(m-2)b=\frac m2+(m-2)b>\frac m2+\frac{m-2}{2}=m-1,
\]
so \(f_{1\perp}(\mu_{0\perp})>m-1\). The unique root \(\mu_{1\perp}\) lies strictly to the left of \(\mu_{0\perp}\).

\end{proof}

\subsection{Ordering of the decoding curvatures}\label{sc:sub9}

For \(\lambda>0\) set \(\mu=\lambda/\sqrt\alpha\) and
\[
\omega_\lambda(c)=\frac{\mu e^{-\mu^2/2}}{m}\,e^{\mu c}\phi(c)^2\Phi(c)^{m-2}.
\]
Let \(A_\lambda=\int\omega_\lambda h^2\), \(B_\lambda=\int\omega_\lambda\sigma\), and \(C_\lambda=\int\omega_\lambda\). The change of scale between \(\omega_\lambda\,dc\) and \(p_{m,\mu}\,dc\) is a positive constant, so \(C_\lambda>0\) and
\[
\frac{A_\lambda}{C_\lambda}=a(\mu),\qquad \frac{B_\lambda}{C_\lambda}=b(\mu).
\]
Define
\[
\Lambda_0=\alpha A_\lambda,\qquad
\Lambda_1=\frac{A_\lambda+(m-2)B_\lambda}{2},\qquad
\Lambda_\perp=\frac{m-1}{2}C_\lambda.
\]
Dividing by \(C_\lambda\) produces exactly \(L_0\), \(L_1\), and \(L_\perp\). Equality and strict inequality of the \(\Lambda\)'s are therefore the same as those of the \(L\)'s.

\begin{theorem}[ordering]\label{sc:theorem9.1}
Fix an integer \(m\ge 4\) and \(\lambda>0\), and let \(\mu=\lambda/\sqrt\alpha\). On the four open intervals determined by Theorem~\ref{sc:theorem8.3},
\begin{align*}
0<\mu<\mu_{1\perp}&\Longrightarrow \Lambda_0<\Lambda_1<\Lambda_\perp,\\
\mu_{1\perp}<\mu<\mu_{0\perp}&\Longrightarrow \Lambda_0<\Lambda_\perp<\Lambda_1,\\
\mu_{0\perp}<\mu<\mu_{01}&\Longrightarrow \Lambda_\perp<\Lambda_0<\Lambda_1,\\
\mu>\mu_{01}&\Longrightarrow \Lambda_\perp<\Lambda_1<\Lambda_0.
\end{align*}
At \(\mu_{1\perp}\) one has \(\Lambda_1=\Lambda_\perp\) and \(\Lambda_0\) strictly smaller than both. At \(\mu_{0\perp}\) one has \(\Lambda_0=\Lambda_\perp\) and \(\Lambda_1\) strictly larger than both. At \(\mu_{01}\) one has \(\Lambda_0=\Lambda_1\) and \(\Lambda_\perp\) strictly smaller than both. The corresponding thresholds in the variable \(\lambda\) are \(\sqrt\alpha\) times the thresholds in \(\mu\).

\end{theorem}
\begin{proof}
Strict increase of \(a+(m-2)b\), of \(a\), and of \(a/b\) converts the three roots into sign changes of \(\Lambda_1-\Lambda_\perp\), of \(\Lambda_0-\Lambda_\perp\), and of \(\Lambda_0-\Lambda_1\), respectively. The order of the roots forbids every other permutation. Lemma~\ref{sc:lemma8.2} is the endpoint comparison: when \(a=m/2\), \(b>1/2\) forces \(\Lambda_\perp<\Lambda_1\), and when \(a=mb\), the same bound forces \(\Lambda_\perp<\Lambda_0\). The relation \(\lambda=\mu\sqrt\alpha\) is the definition \(\mu=\lambda/\sqrt\alpha\).

\end{proof}

\begin{corollary}[the strict normal region]\label{sc:normal-region}
For \(m\ge4\) and \(\lambda>0\), the normal curvature is strictly smaller than both in-plane curvatures if and only if
\(\lambda/\sqrt\alpha>\mu_{0\perp}(m)\).
Moreover,
\[
 \min(\Lambda_0,\Lambda_\perp)<\Lambda_1
\]
for every positive channel parameter.
\end{corollary}
\begin{proof}
The equivalence follows from Theorem~\ref{sc:theorem9.1}, including its endpoint comparisons. The strict inequality is Lemma~\ref{sc:lemma8.2} multiplied by \(C_\lambda>0\).
\end{proof}

\subsection{The triangular case}\label{sc:sub10}

For \(m=3\) the summand \(\mathcal H_0\) is absent. The identity \eqref{sc:identity} becomes \(2a+2b=3+\mu H_3(\mu)\), so
\[
\frac{\Lambda_1}{C_\lambda}=\frac{a+b}{2}=\frac{3+\mu H_3(\mu)}{4},\qquad \frac{\Lambda_\perp}{C_\lambda}=1.
\]
Thus \(\Lambda_1<\Lambda_\perp\), \(\Lambda_1=\Lambda_\perp\), or \(\Lambda_1>\Lambda_\perp\) according as \(\mu H_3(\mu)<1\), \(=1\), or \(>1\).

\begin{theorem}\label{sc:theorem10.1}
The function \(\mu\mapsto\mu H_3(\mu)\) is continuous and strictly increasing from \([0,\infty)\) onto \([0,\infty)\). It has a unique positive root \(\mu_*(3)\) of the equation \(\mu H_3(\mu)=1\). For \(0<\mu<\mu_*\) the in-plane coefficient is strictly smaller; for \(\mu>\mu_*\) the normal coefficient is strictly smaller; at \(\mu_*\) they agree.

\end{theorem}
\begin{proof}
Lemma~\ref{sc:lemma5.1} gives \(H_3'(\mu)=\mathrm{Cov}(h,c)>0\). Hence
\[
\frac{d}{d\mu}\bigl(\mu H_3(\mu)\bigr)=H_3(\mu)+\mu H_3'(\mu)>0
\]
for every \(\mu\ge 0\), the value \(H_3(\mu)\) being positive. The function vanishes at \(0\). Lemma~\ref{sc:lemma8.1} sends it to infinity. The intermediate-value theorem and strict increase give a unique positive root, and the sign chart is the sentence above.

\end{proof}

\section{Proof of the exact normal-shell theorem}\label{ap:normal}

We prove Theorems~\ref{thm:normal-shell} and~\ref{thm:normal-taylor}.
Throughout, $m\ge3$, $n=m-1$, $\alpha=n/m$,
$J=\one\one^{\mathsf T}$, $P=I_m-J/m$, and $G_\Delta=P/\alpha$.
The reference configuration $V=[v_1\ \cdots\ v_m]$ satisfies
$V^{\mathsf T}V=G_\Delta$, and $d\ge m$.
Fix $\lambda>0$ and put $\mu=\lambda/\sqrt\alpha$.
The scalar results in Appendix~\ref{ap:transitions} give the strict normal
regime: $a_m(\mu)>m/2$ if $m\ge4$, and
$a_3(\mu)+b_3(\mu)>2$ if $m=3$.

We use coordinates on the positive square root of the centered Gram matrix.
They give a polynomial parametrization of the extended shell objective;
no differentiable choice of eigenvectors or implicit shell equation is needed.
The objective is extended to a full neighborhood in a fixed finite-dimensional
matrix space before it is restricted to realizable configurations.

\subsection{A smooth extension of the actual decoding deficit}\label{ng:smooth}

For a centered Gaussian vector $\xi$ with covariance $G$, the success
probability of a unit configuration is
\begin{equation}\label{ng:pc}
 \Pc(G;\lambda)=\frac{e^{-\lambda^2/2}}m
                 \E e^{\lambda\max_i\xi_i}.
\end{equation}
This formula is valid when $G$ is singular. If $Z$ is a standard Gaussian
in the signal space, the score vector
$(\langle x_i,Z\rangle)_{i=1}^m$ has exactly covariance $G$.

Let $b$ be an independent standard real Gaussian. The vector
$\sqrt\alpha\,\xi+m^{-1/2}b\one$ has covariance
$R=\alpha G+J/m$. Since $\mu\sqrt\alpha=\lambda$,
\[
 \E_R e^{\mu\max_i X_i}
   =e^{\mu^2/(2m)}\E_G e^{\lambda\max_i\xi_i},
 \qquad
 e^{-\mu^2/2}e^{\mu^2/(2m)}=e^{-\lambda^2/2}.
\]
Consequently the actual deficit is the restriction of
\begin{equation}\label{ng:extended-loss}
 \mathcal F_\lambda(R)=
 \frac{e^{-\mu^2/2}}m
 \left(\E_{I_m}e^{\mu\max_i Z_i}
             -\E_R e^{\mu\max_i X_i}\right):
 \qquad
 \mathcal F_\lambda(\alpha G+J/m)
   =\Pc(G_\Delta;\lambda)-\Pc(G;\lambda).
\end{equation}
The expectation defining this extension makes sense for every positive
definite symmetric $R$ near $I_m$, including matrices for which
$(R-J/m)/\alpha$ is not positive semidefinite.

Here is the regularity justification used below. On a closed covariance
neighborhood satisfying
\begin{equation}\label{ng:covariance-window}
 \tfrac12 I_m\preceq R\preceq2I_m,
\end{equation}
write
\[
 p_R(x)=(2\pi)^{-m/2}(\det R)^{-1/2}
                  \exp(-x^{\mathsf T}R^{-1}x/2).
\]
A derivative of order $k$ in the covariance entries is
$P_{k,R}(x)p_R(x)$, where $P_{k,R}$ is a polynomial of degree at most $2k$
whose coefficients are uniformly bounded on this neighborhood. Hence
\begin{equation}\label{ng:domination}
 |\partial_R^\beta p_R(x)|
 \le K_\beta(1+\norm{x}^{2|\beta|})e^{-\norm{x}^2/4}.
\end{equation}
Multiplication by $e^{\mu\max_i x_i}\le e^{\mu\norm{x}}$ preserves
integrability. Dominated differentiation therefore proves that
$\mathcal F_\lambda$ is smooth on a smaller neighborhood.
In particular it is $C^2$ jointly in all its covariance parameters, and
its third derivatives are locally bounded. All these neighborhoods and
bounds depend only on $m$ and $\lambda$.

For clarity, both bounds in \eqref{ng:covariance-window} are used:
the upper bound supplies the uniform Gaussian decay in
\eqref{ng:domination}. The argument estimates derivatives of the density,
rather than asserting a polynomial bound, uniform in $x$, on a density ratio.

\subsection{The first two covariance variations}\label{ng:variation}

Let $A=A^{\mathsf T}$ have zero diagonal, and set
\[
 S(A)=\sum_{i<j}A_{ij}=\tfrac12\one^{\mathsf T}A\one,
 \qquad N(A)=\norm{A}_F^2,\qquad R_1(A)=\norm{A\one}^2.
\]
We recall the scalar notation
\[
 q(c)=\frac{\phi(c)}{\Phi(c)},\qquad h(c)=c+q(c),\qquad
 \sigma(c)=1-q(c)h(c),\qquad k(c)=\phi(c)^2\Phi(c)^{m-2},
\]
and define the positive weight and its moments by
\begin{align}
 \omega_\lambda(c)&=\frac{\mu e^{-\mu^2/2}}m e^{\mu c}k(c),
       & C_\lambda&=\int_{\R}\omega_\lambda(c)\,dc,\label{ng:weight}\\
 A_\lambda&=\int_{\R}\omega_\lambda(c)h(c)^2\,dc,
       & B_\lambda&=\int_{\R}\omega_\lambda(c)\sigma(c)\,dc.\nonumber
\end{align}
Thus $C_\lambda>0$, $A_\lambda/C_\lambda=a_m(\mu)$, and
$B_\lambda/C_\lambda=b_m(\mu)$.

\begin{lemma}[Full covariance expansion]\label{ng:covariance-expansion}
For symmetric zero-diagonal $A$ near zero,
\begin{align}
 \mathcal F_\lambda(I_m+A)
 &=C_\lambda S(A)+\frac{A_\lambda}{4}N(A)
      -\frac{C_\lambda-B_\lambda}{2}R_1(A)\nonumber\\
 &\quad+\frac{S(A)^2}{2}
            \int_{\R}\omega_\lambda(c)q(c)^2\,dc
       +O(\norm{A}_F^3).
 \label{ng:full-expansion}
\end{align}
The remainder is uniform on a sufficiently small matrix neighborhood.
\end{lemma}

\begin{proof}
At the identity, the density Taylor coefficients are
\begin{align}
 p_{I_m+A}(x)
 =p_{I_m}(x)\bigl[
  1+\tfrac12x^{\mathsf T}Ax+\tfrac14N(A)
   +\tfrac18(x^{\mathsf T}Ax)^2
   -\tfrac12x^{\mathsf T}A^2x\bigr]+\mathcal R_A(x).
 \label{ng:density-expansion}
\end{align}
This follows by expanding $(I_m+A)^{-1}$ and
$-\tfrac12\log\det(I_m+A)$, using $\tr A=0$.
The density differentiation argument above gives, after shrinking the
neighborhood,
\[
 |\mathcal R_A(x)|
 \le K\norm{A}_F^3(1+\norm{x}^6)e^{-\norm{x}^2/4}.
\]
It thus also gives the integrated third-order remainder.

We identify the coefficients by integrating over $(-\infty,c]^m$.
The three one-dimensional truncated moments needed are
\[
 \int_{-\infty}^c\phi=\Phi(c),\qquad
 \int_{-\infty}^c x\phi(x)\,dx=-\phi(c),\qquad
 \int_{-\infty}^c x^2\phi(x)\,dx=\Phi(c)-c\phi(c).
\]
The linear coefficient of $A_{ij}$, $i<j$, is $k(c)$.
For the quadratic coefficients, an edge square $A_{ij}^2$ contributes
$c^2k(c)/2$; a product of distinct edges with one common vertex contributes
$-c q(c)k(c)$; and a product of disjoint edges contributes $q(c)^2k(c)$.
These assertions follow directly from \eqref{ng:density-expansion}:
the corresponding polynomials are
\[
 \tfrac12(1-x_i^2)(1-x_j^2),\qquad
 (x_i^2-1)x_jx_\ell,\qquad x_ix_jx_\ell x_s,
\]
respectively. Each product of distinct edges is counted once.

Put $Q_e=\sum_{i<j}A_{ij}^2=N(A)/2$, and denote the sums of adjacent
and disjoint edge products by $U_e$ and $W_e$. Then
\[
 R_1(A)=2Q_e+2U_e,\qquad
 S(A)^2=Q_e+2U_e+2W_e.
\]
Substitution gives the fixed-$c$ expansion
\begin{align}
 \frac{F_{I_m+A}(c)-\Phi(c)^m}{k(c)}
 =S(A)+\frac{h(c)^2}{4}N(A)
    -\frac{q(c)h(c)}2R_1(A)
    +\frac{q(c)^2}2S(A)^2+O_c(\norm{A}_F^3).
 \label{ng:orthant-expansion}
\end{align}

We justify the subsequent integration, including its sign. Expressing
the expectations through the distribution of the maximum gives
\begin{equation}\label{ng:layer-cake}
 \E_{I_m}e^{\mu\max_i Z_i}-\E_R e^{\mu\max_i X_i}
   =\mu\int_{\R}e^{\mu c}\bigl(F_R(c)-\Phi(c)^m\bigr)\,dc.
\end{equation}
One may obtain it by integration by parts, or by applying Fubini to the
difference of the two laws, which has total mass zero. The same argument
applies to a positive-order covariance derivative of the density:
its integral over all of $\R^m$ is zero. For $c\ge0$, its integral over
$(-\infty,c]^m$ is minus its integral over the complementary set, which is
contained in $\{\norm{x}\ge c\}$. For $c\le0$, the orthant itself lies in
$\{\norm{x}\ge|c|\}$. Formula~\eqref{ng:domination} therefore bounds each
first or second orthant derivative by
$K(1+|c|^N)e^{-c^2/8}$. This remains integrable after multiplication by
$e^{\mu c}$. It validates Fubini and differentiation of
\eqref{ng:layer-cake} through order two.

Multiply the explicitly displayed normalized coefficients in
\eqref{ng:orthant-expansion} by the full weight $\omega_\lambda(c)$,
which includes $k(c)$, and integrate. Since $qh=1-\sigma$, this yields
\eqref{ng:full-expansion}. The bound
$q(c)\le |c|+\sqrt{2/\pi}$ ensures finiteness of the $q^2$ coefficient.
The integrated remainder comes from the smooth density argument, not
from integrating an unspecified fixed-$c$ remainder.
\end{proof}

\subsection{Compression and its exact inverse}\label{ng:compression}

Identify operators on $H=\one^\perp$ with symmetric $m\times m$ matrices
$C$ satisfying $PC=CP=C$. Define the linear map
\begin{equation}\label{ng:uncompress}
 \mathcal U(C)_{ij}=C_{ij}-\tfrac12(C_{ii}+C_{jj}).
\end{equation}
It has zero diagonal, and direct multiplication gives
\begin{equation}\label{ng:inverse}
 P\mathcal U(C)P=C,\qquad
 \mathcal U(PAP)=A
 \quad\text{if }A=A^{\mathsf T},\ \diag A=0.
\end{equation}
For the second identity, write $a=A\one$ and
$s=\one^{\mathsf T}A\one$. The entries of $PAP$ are
$A_{ij}-(a_i+a_j)/m+s/m^2$, and its diagonal entries are
$-2a_i/m+s/m^2$; \eqref{ng:uncompress} cancels the extra terms.
Since a centered matrix has zero total entry sum,
\begin{equation}\label{ng:trace-sum}
 S(\mathcal U(C))=-\frac m2\tr C.
\end{equation}

The shape space used in this appendix is
\[
 \mathcal D_m=\{D=D^{\mathsf T}:PD=DP=D,\ \tr D=0\}.
\]
If $D\in\mathcal D_m$ and $B=\mathcal U(D)$, then $S(B)=0$ and $D=PBP$.
Writing $b=B\one$, we have $\one^{\mathsf T}b=0$ and
\[
 D=B-\frac1m(b\one^{\mathsf T}+\one b^{\mathsf T}).
\]
Expansion of the squared Frobenius norm gives the useful exact identity
\begin{equation}\label{ng:metric-identity}
 \tr(D^2)=\norm{D}_F^2
       =N(B)-\frac2mR_1(B).
\end{equation}

\subsection{The actual Procrustes infimum and a uniform neighborhood}
\label{ng:procrustes}

For unit configurations define
\[
 d_P(X,V)^2=\inf_{U\in O(d)}\norm{X-UV}_F^2.
\]
This is the actual configuration distance. The following calculation
both evaluates the infimum near a simplex and exhibits an attaining
orthogonal map.

Let $G=X^{\mathsf T}X$, and set
\begin{equation}\label{ng:root}
 T=(\alpha PGP)^{1/2},\qquad K=T+(I_m-P).
\end{equation}
The square root is positive semidefinite, $PT=TP=T$, and $T^2=\alpha PGP$.
Assume for the moment that $K$ is invertible. Put
\[
 W=XPK^{-1}.
\]
The regular Gram identity implies $V\one=0$.
The matrix $K$ preserves $H$ and its orthogonal complement; on these
spaces it acts as $T$ and the identity, respectively. Consequently
\[
 W^{\mathsf T}W
 =K^{-1}PGPK^{-1}=\alpha^{-1}P=G_\Delta,
 \qquad WP=W,\qquad XP=WT.
\]
Configurations with the same Gram matrix are related by an orthogonal
map of the ambient space. Indeed, sending the columns of one to those
of the other defines an isometry between their spans: the Gram equality
makes the map well-defined and preserves every inner product.
An orthogonal extension over the complementary subspaces gives an element
of $O(d)$. Thus $W=U_*V$ for some $U_*\in O(d)$.

For any other $Y=UV$, the matrices $Y$ and $W$ are centered and have Gram
matrix $P/\alpha$. Positivity of $T$ gives
\[
 0\le\tr\!\left(T(W-Y)^{\mathsf T}(W-Y)\right)
   =\frac2\alpha\tr T-2\tr(TW^{\mathsf T}Y).
\]
Because $XP=WT$ and $YP=Y$, the last trace equals
$\langle X,Y\rangle_F$. It is at most $\tr T/\alpha$,
with equality at $Y=W$. Since $\norm X_F^2=\norm V_F^2=m$, we obtain
\begin{equation}\label{ng:distance}
 d_P(X,V)^2=2m-\frac2\alpha\tr T.
\end{equation}
In particular the infimum is attained in the situation used here.

We now verify uniformly when this situation applies. If
$\norm{X-UV}_F^2=e$, then, for each pair of indices,
\[
 |G_{ij}-(G_\Delta)_{ij}|
 \le\norm{x_i-Uv_i}+\norm{x_j-Uv_j}\le2\sqrt e.
\]
Taking a sequence approaching the infimum gives
\begin{equation}\label{ng:gram-enclosure}
 \norm{G-G_\Delta}_{\mathrm{op}}
 \le\norm{G-G_\Delta}_F\le2m\,d_P(X,V).
\end{equation}
The bound is independent of $d$ and $V$.

On the cone of positive semidefinite $m\times m$ matrices, the maps
\[
 G\longmapsto(\alpha PGP)^{1/2}+(I_m-P),\qquad
 G\longmapsto\alpha G+J/m
\]
are continuous and take the value $I_m$ at $G_\Delta$.
Invertibility is an open condition. The second map is positive
semidefinite on this cone, so invertibility makes it positive definite.
There is therefore an $\varepsilon_m>0$, in this fixed matrix space,
such that both required conclusions hold when
$\norm{G-G_\Delta}_{\mathrm{op}}<\varepsilon_m$.
Together with \eqref{ng:gram-enclosure}, this supplies a positive
Procrustes neighborhood depending only on $m$. We have used continuity
of the square root on the positive cone, including at the singular
matrix $P$, but have not differentiated that square root.

\subsection{Polynomial square-root coordinates and enclosure}
\label{ng:chart}

Put $\tau=\rho^2$ and
\[
 \kappa=\kappa(\tau)=1-\frac{\tau}{2m}.
\]
If $X$ is on a sufficiently small actual shell $d_P(X,V)^2=\tau$,
\eqref{ng:distance} gives $\tr T=(m-1)\kappa$. Hence
\begin{equation}\label{ng:actual-deviation}
 D=T-\kappa P\in\mathcal D_m.
\end{equation}
Conversely, on the fixed vector space $\mathcal D_m$ define the polynomial
maps
\begin{equation}\label{ng:polynomial-chart}
 C_\tau(D)=(\kappa P+D)^2-P,\qquad
 R_\tau(D)=I_m+\mathcal U(C_\tau(D)),\qquad
 f_\tau(D)=\mathcal F_\lambda(R_\tau(D)).
\end{equation}
For $(\tau,D)$ near $(0,0)$, $R_\tau(D)$ is positive definite and
$f_\tau(D)$ is jointly $C^2$. This definition permits all small
$D\in\mathcal D_m$, irrespective of whether the corresponding Gram
matrix is realizable.

Every actual sufficiently small shell is represented by this chart.
Indeed, for its $D$ in \eqref{ng:actual-deviation},
\[
 C_\tau(D)=T^2-P=P(\alpha G-P)P.
\]
The matrix $\alpha G-P$ has zero diagonal. The inverse identity
\eqref{ng:inverse} gives the exact reconstruction
\begin{equation}\label{ng:reconstruction}
 R_\tau(D)=\alpha G+J/m.
\end{equation}
There is also a quantitative enclosure that uses no modulus of continuity
for the square root. Since the columns of $X$ are unit vectors,
\begin{align*}
 \tr(T^2)
 &=\alpha\tr(PGP)
   =\alpha\left(m-\frac{\norm{\sum_i x_i}^2}{m}\right)
   \le m-1,\\
 \norm D_F^2
 &=\tr(T^2)-\kappa^2(m-1)
   \le(m-1)(1-\kappa^2)\le\tau.
\end{align*}
Thus
\begin{equation}\label{ng:whole-shell-enclosure}
 \norm D_{\mathrm{op}}\le\norm D_F\le\rho
\end{equation}
for every configuration on every sufficiently small positive shell.
The entire feasible shell, not just a chosen smooth path through the
simplex, is covered.

\subsection{Stationarity and the exact shell Hessian}\label{ng:hessian}

Conjugation by a permutation matrix preserves $P$, $\mathcal D_m$,
the map $\mathcal U$, and the chart \eqref{ng:polynomial-chart}.
If $X\sim N(0,R)$ and $Q$ is a permutation matrix, then
$QX\sim N(0,QRQ^{\mathsf T})$, while $\max_i(QX)_i=\max_iX_i$.
Together with the equivariance of the chart, this proves that $f_\tau$
is permutation invariant.
For each sufficiently small fixed $\tau$, its differential at zero is
an invariant linear functional on $\mathcal D_m$. The average of any
$D\in\mathcal D_m$ under all permutations is a matrix $aI_m+bJ$.
Centeredness gives $a+mb=0$, while zero trace gives $a+b=0$.
Since $m>1$, both coefficients vanish. Applying the invariant functional
to this average proves
\begin{equation}\label{ng:stationarity}
 \mathrm d f_\tau(0)[D]=0\qquad(D\in\mathcal D_m).
\end{equation}

We compute the Hessian at $(\tau,D)=(0,0)$.
Fix $D\in\mathcal D_m$ and write $B=\mathcal U(D)$ and
$L=\mathcal U(D^2)$. The polynomial chart gives exactly
\[
 R_0(tD)=I_m+2tB+t^2L.
\]
Here $S(B)=0$ and, by \eqref{ng:trace-sum},
$S(L)=-m\tr(D^2)/2$.
Substitution in \eqref{ng:full-expansion} yields
\begin{equation}\label{ng:Q}
 f_0(tD)=t^2 Q(D)+O(t^3),\qquad
 Q(D)=-\frac{mC_\lambda}{2}\tr(D^2)
      +A_\lambda N(B)-2(C_\lambda-B_\lambda)R_1(B).
\end{equation}
The $S(A)^2$ term of the full expansion is of order $t^4$ here.
It is needed in the general covariance expansion but contributes
nothing to this Hessian. Since $f_0$ is $C^2$,
\begin{equation}\label{ng:actual-hessian}
 \mathrm d^2 f_0(0)[D,D]=2Q(D).
\end{equation}
This follows either by the second derivative chain rule or by comparing
the $t^2$ Taylor coefficient. No remainder symbol is differentiated.

For comparison with the in-plane sectors, write $b=B\one$ and set
\[
 (B_1)_{ii}=0,\qquad
 (B_1)_{ij}=\frac{b_i+b_j}{m-2}\ (i\ne j),\qquad B_0=B-B_1.
\]
Since $\one^{\mathsf T}b=0$, these matrices satisfy
$B_0\one=0$, $B_1\one=b$, and $\langle B_0,B_1\rangle_F=0$.
Direct summation also gives
$\norm{B_1}_F^2=2\norm b^2/(m-2)$.
Define
\[
 E_0=\frac{\norm{B_0}_F^2}{4\alpha},\qquad
 E_1=\frac{\norm b^2}{m-2}.
\]
They are nonnegative and vanish together only when $D=0$. Moreover,
\[
 N(B)=4\alpha E_0+2E_1,\qquad R_1(B)=(m-2)E_1.
\]
For $m=3$, the symmetric zero-diagonal row-zero matrix $B_0$ is zero:
its three row equations force all three off-diagonal entries to vanish.
Thus $E_0=0$ in that case.

With
\[
 \Lambda_0=\alpha A_\lambda,\qquad
 \Lambda_1=\frac{A_\lambda+(m-2)B_\lambda}{2},\qquad
 \Lambda_\perp=\frac{m-1}{2}C_\lambda,
\]
identities \eqref{ng:metric-identity} and \eqref{ng:Q} give
\begin{equation}\label{ng:hessian-gaps}
 Q(D)=4\bigl[(\Lambda_0-\Lambda_\perp)E_0
               +(\Lambda_1-\Lambda_\perp)E_1\bigr],
 \qquad
 \mathrm d^2 f_0(0)[D,D]
   =8\bigl[(\Lambda_0-\Lambda_\perp)E_0
               +(\Lambda_1-\Lambda_\perp)E_1\bigr].
\end{equation}
The factors in this display are part of the coordinate normalization:
the covariance velocity in the present chart is $2B$, and the Hessian
is twice the quadratic Taylor coefficient.

The scalar gap assumptions make \eqref{ng:hessian-gaps} positive definite.
One can see the required quantitative coercivity without any additional
representation-theoretic argument. Formula~\eqref{ng:metric-identity}
reduces \eqref{ng:Q} to
\begin{equation}\label{ng:raw-coercivity}
 Q(D)=\left(A_\lambda-\frac m2C_\lambda\right)N(B)
                       +(2B_\lambda-C_\lambda)R_1(B).
\end{equation}
For $m\ge4$, Appendix~\ref{ap:transitions} gives
$b_m(\mu)>1/2$ for $\mu\ge0$, while the strict normal regime gives
$a_m(\mu)>m/2$. Both coefficients in \eqref{ng:raw-coercivity}
are positive, and $N(B)\ge\norm D_F^2$.
For $m=3$, $R_1(B)=N(B)/2$, so
\[
 Q(D)=(A_\lambda+B_\lambda-2C_\lambda)N(B),
\]
again with a positive coefficient. In either case there is a
$\gamma=\gamma(m,\lambda)>0$ such that
$Q(D)\ge\gamma\norm D_F^2$ for every $D\in\mathcal D_m$.

\subsection{Uniform convexity and the equality orbit}\label{ng:minimum}

Continuity of the Hessian of the joint $C^2$ map $(\tau,D)\mapsto f_\tau(D)$,
together with \eqref{ng:actual-hessian}, gives positive numbers
$\tau_*$, $\delta$, and $\eta$, depending only on $m$ and $\lambda$, such that
\[
 \mathrm d^2f_\tau(D)[H,H]\ge\eta\norm H_F^2
 \quad\text{if }|\tau|<\tau_*,\ \norm D_F<\delta,\
                       H\in\mathcal D_m.
\]
Shrink the neighborhoods if necessary to keep the covariance positive
definite. The ball in $\mathcal D_m$ is convex. Integrating the second
derivative along $t\mapsto tD$, and using \eqref{ng:stationarity}, gives
\begin{equation}\label{ng:strict-minimum}
 f_\tau(D)-f_\tau(0)
 =\int_0^1(1-t)\,\mathrm d^2 f_\tau(tD)[D,D]\,dt
 \ge\frac\eta2\norm D_F^2.
\end{equation}
In particular, $f_\tau(D)=f_\tau(0)$ holds if and only if $D=0$.

Choose $\rho_0>0$ smaller than $\sqrt{2m}$, $\sqrt{\tau_*}$,
$\delta$, and the uniform Procrustes neighborhood from
Subsection~\ref{ng:procrustes}. These choices involve only $m$ and
$\lambda$, before $d$ or $V$ is selected.
By \eqref{ng:whole-shell-enclosure}, every actual unit configuration on
a shell $0<\rho<\rho_0$ satisfies the hypotheses of
\eqref{ng:strict-minimum}. By \eqref{ng:reconstruction} and
\eqref{ng:extended-loss}, its value $f_{\rho^2}(D)$ is its actual
decoding deficit.

It remains to verify attainment and identify every equality case.
Since $\mathcal U(P)=I_m-J$, the point $D=0$ has
\[
 R_{\rho^2}(0)=\kappa^2I_m+(1-\kappa^2)J,\qquad
 G_\kappa=\kappa^2G_\Delta+(1-\kappa^2)J,
 \quad \kappa=1-\frac{\rho^2}{2m}\in(0,1).
\]
The span of $V$ has dimension $m-1$, so $d\ge m$ provides a unit
$z\perp\operatorname{span}\{v_i\}$. Define
\begin{equation}\label{ng:normal-realization}
 N_i=\kappa v_i+\sqrt{1-\kappa^2}\,z.
\end{equation}
These are unit vectors and have Gram matrix $G_\kappa$.
Their centered square root is $T=\kappa P$, so $K$ is invertible and
\eqref{ng:distance} gives
\[
 d_P(N,V)^2=2m(1-\kappa)=\rho^2.
\]
Thus the minimizer of the extended problem is feasible in the original
problem. Its Gram matrix has eigenvalues $\kappa^2/\alpha$ on $H$
and $m(1-\kappa^2)$ on $\operatorname{span}\{\one\}$, so it has rank $m$.

For an arbitrary configuration on this shell, equality in the deficit
bound forces $D=0$ by \eqref{ng:strict-minimum}, hence $G=G_\kappa$ by
reconstruction. The same-Gram isometry argument in
Subsection~\ref{ng:procrustes} then gives precisely the $O(d)$ orbit
of \eqref{ng:normal-realization}. Conversely every configuration on
this orbit has the same Gram matrix and success probability.
Different choices of the unit normal, and permutations of the regular
simplex vertices, produce the same orbit. This proves the equality
assertion and the uniform radius statement of
Theorem~\ref{thm:normal-shell}.

\subsection{The exact value and the sixth-order remainder}\label{ng:taylor}

Let $\xi\sim N(0,G_\Delta)$ and let $b$ be an independent standard real
Gaussian. The vector
$\kappa\xi+\sqrt{1-\kappa^2}\,b\one$ has covariance $G_\kappa$.
For $\kappa\ge0$, its coordinate maximum is
$\kappa\max_i\xi_i+\sqrt{1-\kappa^2}\,b$.
Using the Gaussian moment-generating function in \eqref{ng:pc} yields
\begin{equation}\label{ng:exact-value}
 \Pc(G_\kappa;\lambda)
   =\frac{e^{-\lambda^2/2}}m
      e^{\lambda^2(1-\kappa^2)/2}
      \E e^{\lambda\kappa\max_i\xi_i}
   =\Pc(G_\Delta;\lambda\kappa).
\end{equation}
This proves the exact optimal value, with all prefactors included.

Write $p(s)=\Pc(G_\Delta;s)$ using the expectation formula as a real
extension in $s$. This extension is smooth. Indeed, if
$M=\max_i\xi_i$, then each coordinate $\xi_i$ is Gaussian and
$|M|\le\max_i|\xi_i|$. Therefore
$\E(|M|^k e^{L|M|})<\infty$ for every fixed $k$ and $L$.
On compact $s$-intervals these moments dominate all the derivatives of
$e^{-s^2/2}\E e^{sM}/m$. In particular $p$ is $C^3$ near $\lambda$,
with bounds depending only on $m$ and $\lambda$.

Choose $u>0$ and $M_3>0$ such that
$|p'''(s)|\le M_3$ whenever $|s-\lambda|\le u$.
For $h=\lambda\rho^2/(2m)$ small enough, Taylor's theorem gives
\[
 \left|p(\lambda)-p(\lambda-h)
       -p'(\lambda)h+\frac12p''(\lambda)h^2\right|
       \le\frac{M_3}{6}|h|^3.
\]
Together with \eqref{ng:exact-value}, this proves
\begin{equation}\label{ng:O6}
 \delta_{\min}(\rho)
 =\frac{\lambda p'(\lambda)}{2m}\rho^2
    -\frac{\lambda^2p''(\lambda)}{8m^2}\rho^4+O(\rho^6).
\end{equation}
More explicitly, after decreasing $\rho_0$ if necessary there is a
$C>0$, depending only on $m$ and $\lambda$, such that
\begin{equation}\label{ng:quantified-remainder}
 \left|\delta_{\min}(\rho)
       -\frac{\lambda p'(\lambda)}{2m}\rho^2
       +\frac{\lambda^2p''(\lambda)}{8m^2}\rho^4\right|
 \le C|\rho|^6,\qquad 0<\rho<\rho_0.
\end{equation}
For example the Taylor argument permits $C=M_3\lambda^3/(48m^3)$
on a sufficiently small interval. The same estimate holds on a
two-sided interval for the real extension
$p(\lambda)-p(\lambda(1-\rho^2/(2m)))$.

That exact extension is even in $\rho$, so every odd derivative at zero
vanishes. In particular there is no cubic term; this follows from the
exact value, not from dropping an uncomputed cubic remainder.
The kernel-checked Taylor argument proves the same quantified
$O(\rho^6)$ statement by applying the order-three Taylor little-$o$
theorem, adding back its cubic coefficient to obtain
$O(h^3)$ after the quadratic polynomial, and composing with
$h=-\lambda\rho^2/(2m)$. It thus verifies the full sixth-order bound,
in addition to the parity identity.

The proof uses strict positivity of the existing in-plane curvature
gaps. At a threshold where one such gap vanishes,
\eqref{ng:strict-minimum} is not supplied by this argument.
No assertion about that degenerate case, or an effective numerical
lower bound for $\rho_0$, is part of the theorem.

\section{Lean formalization and reproducibility}\label{ap:lean}

\subsection{Scope of the certificate}\label{lv:scope}
The supplementary Lean development verifies the global quantitative
orthant comparison with its explicit constants, the separate centroid
penalty, and their consequences for decoding, Procrustes recovery, rank
deficiency, unused energy and mean width.  This extends the earlier
certificate for the ordered channel-parameter crossings, the exact
strict-normal minimizer on sufficiently small Procrustes shells, and the
quantified $O(\rho^6)$ Taylor remainder.  All statements use the actual
Gaussian measures, integrals and geometric distance.  The global
comparison does not assume a witness bound, a product deficit or a
desired CDF inequality: these are proved in the dependency chain,
including for singular covariances.

The global results hold for $m\ge2$, with $\lambda>0$ where a positive
decoding coefficient is required.  This does not relax the separate
hypotheses of the local results: those retain $m\ge3$, the stated strict
curvature gap, and a sufficiently small shell radius.  The certificate
does not cover numerical diagnostics, learning experiments or historical
priority claims.  It does not identify finite-radius minimizers in the
tangential regimes or on a boundary of equal curvatures.  The local
radius and remainder constant are proved to exist; no effective numerical
value is certified.

The final library entry is \path{SimplexFull/VerifiedGlobal.lean}.
Its dependency closure contains 216 source modules, including
52 vendored foundation modules.  A separate audit imports this entry
and prints the axioms of 91 selected public declarations.
These counts describe the supplied artifact, rather than a claim that
every sentence of the manuscript has a formal counterpart.  The formal
and written proofs may use different intermediate lemmas.
The original 94-module local development is retained byte for byte.

\subsection{Objects and normalization}\label{lv:objects}
The Gaussian density and distribution function are concrete definitions:
\path{SimplexFull.phi} is Mathlib's real standard Gaussian density, and
\path{SimplexFull.Phi} is its Lebesgue integral over $(-\infty,c]$.
The definitions of $q,h,\sigma$ therefore have their usual analytic
meaning.  For $m\geq3$, the normalized moments in the crossing theorems are
defined using
\[
 w_{m,\mu}(c)=e^{\mu c}\phi(c)^2\Phi(c)^{m-2},
 \qquad
 Z_m(\mu)=\int_{\mathbb R}w_{m,\mu}(c)\,dc,
\]
and
\[
 a_m(\mu)=\frac{\int h(c)^2w_{m,\mu}(c)\,dc}{Z_m(\mu)},\qquad
 b_m(\mu)=\frac{\int \sigma(c)w_{m,\mu}(c)\,dc}{Z_m(\mu)},\qquad
 H_m(\mu)=\frac{\int h(c)w_{m,\mu}(c)\,dc}{Z_m(\mu)}.
\]
These are respectively \path{SimplexFull.a}, \path{SimplexFull.b}, and
\path{SimplexFull.H}; their normalizer is proved positive.  The parameter
in these integrals is $\mu=\lambda/\sqrt{\alpha}$, where
$\alpha=(m-1)/m$.  Thus a crossing expressed in $\mu$ corresponds to the
physical signal amplitude $\lambda=\sqrt{\alpha}\,\mu$.  The variables
$c$, $\mu$, and $\lambda$ are not interchangeable.  For $\lambda>0$,
\[
 a_m(\mu)=A_\lambda/C_\lambda,\qquad
 b_m(\mu)=B_\lambda/C_\lambda,\qquad
 C_\lambda=\frac{\mu e^{-\mu^2/2}}{m}Z_m(\mu)>0.
\]
The value $\mu=0$ used in the normalized crossing arguments is an endpoint
of the moment family, not a division by the zero-amplitude value of
$C_\lambda$.

The objective \path{SimplexFull.decodingSuccess} is defined by
\[
 \Pc(G;\lambda)=\frac{e^{-\lambda^2/2}}{m}
 \int_{\mathbb R^m}e^{\lambda\max_i x_i}\,dN(0,G)(x).
\]
Here $N(0,G)$ is Mathlib's actual
\path{ProbabilityTheory.multivariateGaussian}, including the singular
simplex covariance.  Gaussian density identities for positive definite
covariances and the common Gaussian shift identity connect the local
analytic calculations to this definition.  The extended development also identifies this expression with the
operational maximum-likelihood decoder for unit signals.  For signals
of unequal energy it instead uses the class likelihood
$\exp(\lambda\langle x_i,z\rangle-\lambda^2\|x_i\|^2/2)$ and the
maximum-likelihood/Bayes-value identity.  The individual energy correction
is retained when proving the unused-energy result.

In the namespace \path{SimplexFull.Geometry},
\path{simplexSuccess m} is exactly the function
$p_m(\lambda)=\Pc(G_\Delta;\lambda)$.  The symbols $p_m'$ and $p_m''$
in the Taylor statement are represented by \path{deriv} and
\path{deriv (deriv (simplexSuccess m))}, rather than by independent
coefficient variables.  The theorem
\path{SimplexFull.contDiff_simplexSuccess} supplies genuine $C^3$
regularity.  Its proof establishes all-real exponential integrability
of the coordinate maximum and applies Mathlib's moment-generating-function
analyticity theorem.  In fact, the formal regularity result is stronger
than the order needed here.

For configurations, \path{Geometry.procrustesSq X V} is defined as the
real infimum
\[
 \inf_{U\in O(d)}\sum_{i=1}^m\|X_i-U V_i\|^2.
\]
The index type consists of linear isometric automorphisms of the whole
Euclidean space.  Accordingly, \path{Geometry.orthoOrbit X N} expresses
membership in the actual $O(d)$ orbit, including reflections.  The
distance is not replaced by a chosen alignment or an auxiliary matrix
norm.  The normal configuration is proved realizable when $d\geq m$,
and its squared Procrustes distance is proved to equal $\rho^2$.

For a finite configuration $X$, the support of its actual convex hull is
defined as the supremum of $\langle g,y\rangle$ over
$y\in\operatorname{conv}(\operatorname{range}X)$.  The development proves
that this supremum is attained and equals the finite maximum of the
scores.  The normalized sphere measure is the normalization of
Mathlib's \path{volume.toSphere}; it is not defined through a Gaussian
ratio.  The polar-coordinate measure-preserving map, the radial Gaussian
density, and positive homogeneity identify twice its support average (equivalently, the average of
$h_K(u)+h_K(-u)$) with the
Gaussian expression $2\E\max_i\langle g,X_i\rangle/\E\|g\|$.
The denominator is always the Gaussian radial mean in the original
ambient dimension $d$, including when an auxiliary lift uses $d+m$
coordinates.

\subsection{Correspondence with the new theorems}\label{lv:map}
Table~\ref{lv:table} lists the principal entry declarations.  Module paths
are relative to \path{SimplexFull}; all declarations shown below are
fully qualified.  The package's proof map and source index give the
remaining dependency information.

\begingroup
\small
\renewcommand{\arraystretch}{1.18}
\begin{longtable}{@{}>{\raggedright\arraybackslash}p{0.25\textwidth}p{0.70\textwidth}@{}}
\caption{Principal mathematical statements and their Lean entry points.}
\label{lv:table}\\
\toprule
Statement & Module and public declaration\\
\midrule
\endfirsthead
\toprule
Statement & Module and public declaration\\
\midrule
\endhead
Theorem~\ref{thm:global}
&
\path{GlobalComparison.lean}:
\path{SimplexFull.GlobalComparison.orthant_sqrt_comparison},
\path{orthant_linear_comparison}, and
\path{quantitative_lowerOrthant} in the same namespace.
These use the exact $K,\kappa$ and all admissible PSD covariances.
\\
Quantitative rectangle and centroid bounds
&
\path{GlobalRectangleBound.lean}:
\path{SimplexFull.GlobalRectangle.symmetricRectangle_gap};
\path{GlobalCentroidConfig.lean}:
\path{SimplexFull.GlobalCentroid.centroid_penalty}.
\\
Theorem~\ref{thm:code} and Corollary~\ref{cor:operational}
&
\path{GlobalApplicationsCore.lean} and
\path{GlobalApplicationsDecoder.lean}:
\path{SimplexFull.GlobalApplication.decodingSuccess_gram_penalty},
\path{decodingSuccess_procrustes_penalty}, and
\path{decodingSuccess_rank_deficiency} in that namespace.
The centroid and operational unused-energy conclusions are
\path{SimplexFull.GlobalTail.decodingSuccess_centroid_penalty} and
\path{SimplexFull.GlobalTail.ball_decoder_unused_energy}.
\\
Theorem~\ref{thm:width}
&
\path{GlobalSphereGaussianWidth.lean}:
\path{SimplexFull.SphereGaussian.ball_mean_width}.
The traditional width-function integral is identified by
\path{GlobalSphereGaussianSymmetry.lean}:
\path{SimplexFull.SphereGaussian.widthIntegral_convexHull_eq_meanWidth}.
The Gaussian representation is
\path{SimplexFull.GlobalApplication.ball_mean_width}.
\\
Theorem~\ref{thm:ordered-switches}
&
\path{GaussianCrossings.lean}:
\path{SimplexFull.gaussian_ordered_crossings}.
The strict exclusion of the intermediate coefficient from the minimum is
\path{SimplexFull.gaussian_middle_coefficient_not_min}.
The complete chart and its positive-scaling form are in
\path{CrossingSignChart.lean}:
\path{SimplexFull.gaussian_ordered_crossings_sign_chart} and
\path{SimplexFull.gaussian_ordered_crossings_sign_chart_scaled}.
\\
Theorem~\ref{thm:triangular-switch}
&
\path{GaussianCrossings.lean}:
\path{SimplexFull.gaussian_triangular_unique_crossing}.
The equivalent equation $\mu H_3(\mu)=1$ is established by
\path{SimplexFull.gaussian_triangular_coefficient_eq_iff}.
\\
Theorem~\ref{thm:normal-shell}
&
\path{UniformNormalShell.lean}:
\path{SimplexFull.Geometry.gaussian_uniform_exact_normal_shell_m4}
and
\path{SimplexFull.Geometry.gaussian_uniform_exact_normal_shell_m3}.
The corresponding threshold versions are
\path{SimplexFull.Geometry.gaussian_normal_threshold_uniform_shell_m4}
and
\path{SimplexFull.Geometry.gaussian_normal_threshold_uniform_shell_m3}.
\\
Theorem~\ref{thm:normal-taylor}
&
\path{GaussianTaylor.lean}:
\path{SimplexFull.Geometry.normalTaylor_bound},
\path{SimplexFull.Geometry.normalTaylor_isBigO},
\path{SimplexFull.Geometry.gaussian_uniform_taylor_shell_m4},
and
\path{SimplexFull.Geometry.gaussian_uniform_taylor_shell_m3}.
\\
\bottomrule
\end{longtable}
\endgroup

For $m\geq4$, the crossing declaration asserts
$0<t_1<t_0<t_{01}$ and uniqueness on $[0,\infty)$ of the equations
\[
 a_m(t_1)+(m-2)b_m(t_1)=m-1,\qquad
 a_m(t_0)=m/2,\qquad
 a_m(t_{01})=m b_m(t_{01}).
\]
The declaration \path{gaussian_ordered_crossings_sign_chart} packages
these roots and their uniqueness together with all four open-interval
comparisons and all three endpoint comparisons for $L_0,L_1,L_\perp$.
Its \path{gaussian_ordered_crossings_sign_chart_scaled} companion gives the
same chart after multiplication by any positive common factor, with the
roots independent of that factor.  For every $\lambda>0$ one may therefore
use $C_\lambda>0$ and $\Lambda_j=C_\lambda L_j$ to obtain the physical
coefficient chart in Theorem~\ref{thm:ordered-switches}.  No additional
analytic hypothesis is required.  The original crossing declaration
instantiates the abstract comparison lemmas with the actual Gaussian moments.
The proofs include the strict benchmark $b_4(0)>1/2$,
monotonicity, and the integral balance identity
$2a_m+(m-1)b_m=m+\mu H_m$.
The case $m=3$ is handled separately: its crossing is
$a_3+b_3=2$, and no nonzero $\mathcal H_0$ sector is introduced.

The geometric conclusions are packaged as
\path{SimplexFull.Geometry.UniformExactNormalShell} and
\path{SimplexFull.Geometry.UniformTaylorNormalShell}.
For fixed $m$ and $\lambda$ in the strict-normal regime, the latter has
the quantifier order
\[
 \exists C>0\ \exists\rho_0>0\
 \forall d\geq m\ \forall V\text{ with Gram }G_\Delta\
 \exists z\
 \forall\,0<\rho<\rho_0\
 \forall X\text{ on the feasible shell}.
\]
The normal vector $z$ is unit and orthogonal to every $V_i$.
In particular, both $C$ and $\rho_0$ are independent of the ambient
dimension, reference embedding, and competitor.  The chosen $z$ is
independent of $\rho$.  The final statements contain feasibility of
$N_i=\kappa V_i+\sqrt{1-\kappa^2}\,z$, the exact optimal value, a comparison
with every feasible $X$, and equality if and only if
\path{Geometry.orthoOrbit X N}.  The strict-normal hypothesis is
$a_m(\mu)>m/2$ for $m\geq4$, and $a_3(\mu)+b_3(\mu)>2$ for $m=3$.

The Taylor remainder is the actual function
\[
 \mathcal R_{m,\lambda}(\rho)
 =p_m(\lambda)-p_m\!\left(\lambda\left(1-\frac{\rho^2}{2m}\right)\right)
 -\left[
 \frac{\lambda p_m'(\lambda)}{2m}\rho^2
 -\frac{\lambda^2p_m''(\lambda)}{8m^2}\rho^4
 \right].
\]
The declaration \path{normalTaylor_bound} proves
$|\mathcal R_{m,\lambda}(\rho)|\leq C|\rho|^6$ for $|\rho|<r$,
with $C,r>0$.  The companion \path{normalTaylor_isBigO} states the
corresponding two-sided big-$O$ assertion at zero.  The generic proof
uses a third-order Taylor expansion, obtains an $O(h^3)$ quadratic
remainder, and substitutes $h=-\lambda\rho^2/(2m)$.  Thus it proves
$O(\rho^6)$, rather than only the $o(\rho^4)$ conclusion available
from $C^2$ regularity.

The stronger regularity declaration
\path{SimplexFull.contDiff_simplexSuccess_top} also verifies the
analyticity assertion.  In the pinned Mathlib version the order type of
\path{ContDiff} is $\operatorname{WithTop}(\mathbb N_\infty)$:
its outer $\top$ is the analytic order $\omega$, distinct from the
smooth order $\infty$.  Applying \path{ContDiffAt.analyticAt} therefore
gives \path{AnalyticAt} at every real argument.  A supplementary
kernel-checked statement spells out this consequence for the actual
multivariate Gaussian probability integral.

The declaration
\path{SimplexFull.Geometry.normalDeficit_odd_iteratedDeriv_zero}
in \path{TaylorParity.lean} proves that every odd derivative of the exact
deficit vanishes at zero.  Finally,
\path{SimplexFull.Geometry.simplex_quadratic_coefficient_eq_normal_curvature}
in \path{TaylorCoefficient.lean} identifies the leading coefficient with
the actual normal curvature:
\[
 \frac{\lambda p_m'(\lambda)}{2m}
 =\frac{m-1}{2}\,C_\lambda.
\]
This identity includes the conversion between $\lambda$ and $\mu$.


\subsection{Global proof chain and assumptions}\label{lv:global}
The scalar function $\Psi$ is defined from the actual inverse of
$h(s)=s+q(s)$ on the positive axis.  Its monotonicity, inverse range,
strict concavity, and derivatives are proved.  The curvature constant
$M_0$ is the attained maximum of the actual continuous curvature
function on the manuscript's compact interval.  In particular,
the lower Taylor bound used in the centroid argument is not supplied
as an external smoothness or remainder assumption.

The quantitative rectangle proof starts from the true bivariate
conditional Gaussian law.  Differentiation under the compact integral
and the boundary derivative give the Plackett formula.  Integration
on $[0,1/2]$, monotonicity and evenness yield the stated
$\ell^2e^{-4T^2/3}/(4\pi)$ coefficient.  At correlation $+1$ or $-1$
the residual Gaussian is a Dirac measure and the rectangle mass is
the smaller marginal slab mass; no singular density is substituted.
The actual Gaussian slab peeling argument and a maximizing
off-diagonal pair then give the factor $m(m-1)$ and the exponent $m-2$.

For Theorem~\ref{thm:global}, existence of actual adaptive witnesses,
their compact coordinate bounds, the sum--difference Gaussian law
and Fubini's theorem lead to the quantitative doubled-product deficit.
The product of the normalized self-convolution factors has Gaussian
integral at least one, including
for PSD singular covariance by continuous regularization.
The fixed-factor rectangle estimate on the central box gives precisely
$K$, with exponent $2m-2$.  The correlation-matrix bound
$\|R-I\|_F^2\le m(m-1)$ gives the manuscript's $\kappa$ by
rationalizing the square root.  The identity covariance is handled
directly.  The centroid proof uses the actual PSD variational value,
the verified lower Taylor bound for $\Psi$, and the exact $\tau,\eta$.

The integration bridge proves
\[
 \E\max Z-\E\max X
 =\int_{\mathbb R}\bigl(F_X(c)-F_Z(c)\bigr)\,dc
\]
from a centered signed-tail kernel and Fubini, with integrability
proved.  Applied to the exponential of the maximum it retains the
exact window factor $e^\mu-1$.  Gaussian normalization at
$R=\alpha G+J/m$ supplies the factors $\alpha^2$ and $\sqrt\alpha$.
Consequently the final entries use
\[
 c_g=W\kappa\alpha^2,\quad c_s=W\eta,\quad
 c_P=\left(\frac{\alpha}{c_g}+\frac{1}{mc_s}\right)^{-1},
 \qquad b_g=\kappa\alpha^{3/2},\quad b_s=\eta/\sqrt\alpha.
\]
None of those entries retains a unit-comparison or target-CDF hypothesis.

The geometric inequality is proved for vectors of arbitrary lengths
and for the actual infimum over $O(d)$; a polar-alignment witness also
proves attainment.  The rank corollary assumes the actual matrix rank
bound, with its bound independent of any realization dimension.
For ball configurations the formal proof uses mutually orthogonal
extra coordinates.  Adding the resulting independent centered Gaussian
noise cannot decrease the expected maximum; the corresponding likelihood
argument proves the operational decoding comparison for nonunit signals.
This orthogonal lift is an alternative to the shared-coordinate lift
in the written unused-energy proof.

The foundation sources from the weak-simplex and simplex-uniqueness
developments are supplied with exact upstream commits, per-file
provenance and their original licenses and attributions in
\path{VENDOR_PROVENANCE.json} and \path{licenses/}.  Their imported
closure is included in the same fresh trust-level-zero build as the
new modules, rather than accepted through inherited project binaries.

\subsection{Build and axiom audit}\label{lv:certificate}
The released archive is
\path{concrete-simplex-lean-global-20261106.zip}.
Its SHA-256 digest is
\begin{quote}\small
\path{4e8e24831e8c56ab07523aed22c89c4ba6854726accccff2aa32c0587133896b}.
\end{quote}
The toolchain is Lean 4.31.0, pinned by
\path{leanprover/lean4:v4.31.0}; the Mathlib revision is
\begin{quote}\small
\path{fabf563a7c95a166b8d7b6efca11c8b4dc9d911f}.
\end{quote}
The successful recorded run is \path{20261106-180522-470580}.
Its \path{evidence/verification.json} records all 216 source-module
commands and exit codes, source and compiled-file hashes, the successful
combined import, and the final axiom audit.  All source hashes remained
unchanged through verification.  The manuscript bundle also supplies
the independent Grok Build and agent cross-reviews.

The verifier copies the three project-source namespaces into a fresh
snapshot and creates an initially empty output directory.  It compiles
their explicit dependency closure in dependency order.  All inherited
project build paths, including paths containing only a bare root module,
are excluded from \path{LEAN_PATH}; the pinned Mathlib cache is supplied
by Lake.  Every project source is compiled with Lean's trust level zero:
\begin{verbatim}
lean -j 2 -t 0 -R <snapshot/src> \
  -o <fresh/build/Module.olean> \
  <snapshot/src/Module.lean>
\end{verbatim}
The final audit imports \path{SimplexFull.VerifiedGlobal} and invokes
\path{#print axioms} for 91 selected declarations.
Their dependency reports contain only
\path{propext}, \path{Classical.choice}, and \path{Quot.sound}.
Acceptance requires exactly one recognized report per requested
declaration.  Missing, duplicate, unexpected or nonstandard reports
cause verification to fail.

After removing comments and strings, source scanning rejects
\path{sorry}, \path{admit}, project \path{axiom} declarations,
\path{unsafe}, \path{native_decide} and \path{debug.skipKernelTC}.
Import parsing rejects unknown project roots and unsupported header
layouts instead of silently omitting them.  Positive and negative
regression tests cover the parser, binary-path isolation and axiom
reports.  These checks accompany, rather than replace, kernel checking.

\subsection{Reproducing the verification}\label{lv:reproduce}
After extracting the archive, install the pinned Lean toolchain with
\path{elan} and make Python available.  From the package directory run
\begin{verbatim}
lake update
lake exe cache get
lake env python reproduce.py
\end{verbatim}
The Mathlib revision is fixed in \path{lakefile.toml}.
The verifier uses \path{SimplexFull.VerifiedGlobal} and
\path{AUDIT_THEOREMS_GLOBAL.json}, creates a fresh timestamped build,
and returns a nonzero status on any source, compilation or axiom-audit
failure.  The package supplies \path{PROOF_MAP_GLOBAL.md},
\path{SOURCE_INDEX.json}, \path{SHA256SUMS.json},
the individual compiler logs, upstream provenance and independent
cross-reviews.  The source, compiler output and
\path{evidence/verification.json} are the reproducibility evidence.


\begin{thebibliography}{99}
\bibitem{Mulgund} A. Mulgund,
\emph{Stochastic Domination of Gaussian Maxima: A Resolution of the Weak
Simplex Conjecture}, arXiv:2607.14087. The non-strict product and adaptive
tilting inputs used here are in v2, 20 July 2026, Theorem 4.1 and Proposition 5.2.
\url{https://arxiv.org/abs/2607.14087v2}.
\bibitem{Su} M. Su, K. Yang, H. Xu, and Chih-Lin I,
\emph{The Equality Cases of the Weak Simplex Conjecture},
arXiv:2608.19093v1, 19 August 2026.
\url{https://arxiv.org/abs/2608.19093}.
\bibitem{Royen} T. Royen,
\emph{A simple proof of the Gaussian correlation conjecture extended to
some multivariate gamma distributions}, Far East Journal of Theoretical
Statistics 48 (2014), 139--145. \url{https://arxiv.org/abs/1408.1028}.
\bibitem{Litvak} A. E. Litvak,
\emph{Around the simplex mean width conjecture}, in Analytic Aspects of
Convexity, Springer INdAM Series 25 (2018), 73--84.
\url{https://www.math.ualberta.ca/~alexandr/OrganizedPapers/AL-simp-wv.pdf}.
\bibitem{BFH} K. J. B\"or\"oczky, F. Fodor, and D. Hug,
\emph{Strengthened inequalities for the mean width and the \(\ell\)-norm},
Journal of the London Mathematical Society (2) 104 (2021), no. 1, 233--268.
\url{https://doi.org/10.1112/jlms.12429}.
\bibitem{MulgundFull} A. Mulgund,
\emph{Stochastic Domination of Gaussian Maxima by the Regular Simplex},
arXiv:2609.28452v2, 27 September 2026.
\url{https://arxiv.org/abs/2609.28452v2}.
\bibitem{MulgundLatest} A. Mulgund,
\emph{Stochastic Domination of Gaussian Maxima: A Resolution of the Weak
Simplex Conjecture}, arXiv:2607.14087v4, 21 August 2026.
\url{https://arxiv.org/abs/2607.14087v4}.
\bibitem{Balakrishnan} A. V. Balakrishnan,
\emph{A contribution to the sphere-packing problem of communication theory},
Journal of Mathematical Analysis and Applications 3 (1961), 485--506.
\url{https://doi.org/10.1016/0022-247X(61)90072-5}.
\bibitem{LandauSlepian} H. J. Landau and D. Slepian,
\emph{On the optimality of the regular simplex code},
Bell System Technical Journal 45 (1966), 1247--1272.
\url{https://doi.org/10.1002/j.1538-7305.1966.tb01696.x}.

\bibitem{DNS} A. De, S. Nadimpalli, and R. A. Servedio,
\emph{Quantitative Correlation Inequalities via Semigroup Interpolation},
12th Innovations in Theoretical Computer Science Conference (ITCS 2021),
LIPIcs 185, Article 69.
\url{https://doi.org/10.4230/LIPIcs.ITCS.2021.69}.
\bibitem{FGS} O. Friedland, O. Gu\'edon and F. Souli,
\emph{On partition functions of Gaussian random variables},
arXiv:2609.38904v1, 30 September 2026.
\url{https://arxiv.org/abs/2609.38904v1}.
\bibitem{Weber} C. L. Weber,
\emph{Signal design for space communication channels, part I},
USCEE-129 / NASA-CR-78972, 1 February 1965.
\url{https://ntrs.nasa.gov/citations/19660030423}.
\bibitem{OSheaHoydis} T. O'Shea and J. Hoydis,
\emph{An Introduction to Deep Learning for the Physical Layer},
IEEE Transactions on Cognitive Communications and Networking 3 (2017),
563--575. \url{https://arxiv.org/abs/1702.00832v2}.
\end{thebibliography}
\end{document}
